\documentclass[10pt,a4paper]{article}
\usepackage{amsmath,amssymb,amsthm,graphicx}
\usepackage{amsmath,amssymb,amsthm,algorithmicx,algpseudocode}
\usepackage[no-math]{fontspec}
\newenvironment{authorclaim}[1]{\par\medskip\noindent\textbf{#1.}\ \ignorespaces}{\par\medskip}
\allowdisplaybreaks
\title{Multiparty NOF search lower bounds for bit pigeonhole and semantic proof consequences}
\author{Dmitry Itsykson, Artur Riazanov}
\date{doi:10.4230/LIPIcs.CCC.2021.3; CC BY 4.0}
\begin{document}
\maketitle
\noindent\textit{Editorial scope.} Counted once: full bit-PHP multiparty NOF search and semantic consequences. All load-bearing author reasons, probability argument, symmetry/bijection proofs, induction cases, reductions and semantic substitution continuations retained in editable TeX. Non-load-bearing background connective prose may be condensed. Original printed conventions, including Definition23/Claim27 indices and Section6.3 Q polynomial, retained without mathematical repair; formal PDF is the original source.\par\medskip
\noindent\textit{Allowed standard prerequisites (editorial).}\begin{itemize}
\item Randomized communication complexity, unique/set disjointness lower bounds and biased random-walk hitting time. Formal Section2 Theorems7–8 and AppendixA Theorem45, with original full33-entry bibliography. Standard imported results retained as cited prerequisites, not recursively reproved.
\end{itemize}\medskip
% Faithful typesetting transcription of necessary formal Section2 definitions.
\section*{2 Preliminaries}
\paragraph{Notations.} We use the following notation: $[n]=\{1,2,\ldots,n\}$. Let $S^{n\times m}$ denote the set of matrices of size $n\times m$ with elements from $S$. We denote by $0_{n\times m}$ the zero matrix of size $n\times m$ and by $1_{n\times m}$ the matrix of the same size containing only ones. For square matrices $A_1,\ldots,A_k$ we denote a diagonal block matrix with blocks $A_1,\ldots,A_k$ by $\operatorname{diag}(A_1,\ldots,A_k)$. For $x\in\{0,1,\ldots,2^k-1\}$ we denote a vector $(a_0,\ldots,a_{k-1})\in\{0,1\}^k$ such that $x=\sum_{i=0}^{k-1}a_i2^i$ by $\operatorname{bin}_k(x)$, i.e. $(a_0,\ldots,a_{k-1})$ is the reversed binary representation of $x$. For vectors $v_1,\ldots,v_n$ from a vector space over a field $\mathbb F$ we denote their linear span by $\operatorname{Span}(v_1,\ldots,v_n)$. We use coordinate-wise comparison of strings from $\{0,1\}^n$, i.e. for $x,y\in\{0,1\}^n$ we write $x\le y$ iff $x_i\le y_i$ for each $i\in[n]$. We denote the set of variables of a CNF-formula $\phi$ by $\operatorname{Vars}(\phi)$.
\paragraph{Communication complexity.} We briefly recall some notions of communication complexity. For formal definition and details we refer to [24].
In the classic two-party randomized communication protocol with public randomness, Alice and Bob cooperate to compute a relation $Q\subseteq X\times Y\times Z$: Alice has an input $x\in X$ and Bob has an input $y\in Y$, their goal is to compute $z\in Z$ such that $(x,y,z)\in Q$. We assume that Alice and Bob have access to an arbitrary large random string of bits that is common for Alice and Bob. Let for every $x\in X$ and $y\in Y$, $R^\delta_{\mathrm{pub}}(Q,x,y)$ denote the minimal number of bits Alice and Bob need to transmit between each other so they both find a $z\in Z$ such that $(x,y,z)\in Q$ with probability at least $1-\delta$ taken over the values of the common random string. And $R^\delta_{\mathrm{pub}}(Q):=\max_{x\in X,y\in Y}R^\delta_{\mathrm{pub}}(Q,x,y)$.
We also consider multiparty communication protocols in the number on forehead (NOF) model that extends two-party protocols for an arbitrary number of parties. In this setting $k$ parties cooperate to compute a relation $Q\subseteq X_1\times X_2\times\cdots\times X_k\times Y$. The $i$th party has $x_i\in X_i$ written on their forehead so they know all $x_j$ for $j\ne i$, their goal is to compute $y\in Y$ such that $(x_1,x_2,\ldots,x_k,y)\in Q$. The parties communicate by taking turns broadcasting messages to all other parties until all parties learn the value of $y\in Y$ such that $(x_1,\ldots,x_k,y)\in Q$. In this model we also assume that all parties have access to a common random string of bits. Let $R^\delta_{\mathrm{pub}}(Q,x_1,\ldots,x_k)$ for $x_1\in X_1,\ldots,x_k\in X_k$ denote the minimal total number of bits transmitted until each party learns $y\in Y$ such that $(x_1,\ldots,x_k,y)\in Q$ with probability at least $1-\delta$ taken over the set of values of the random string of bits. Also, let $R^\delta_{\mathrm{pub}}(Q):=\max_{x_1\in X_1,\ldots,x_k\in X_k}R^\delta_{\mathrm{pub}}(Q,x_1,\ldots,x_k)$.
Let $f$ be a function from $X_1\times X_2\times\cdots\times X_k\to Y$. Then $R^\delta_{\mathrm{pub}}(f)$ denotes $R^\delta_{\mathrm{pub}}(Q_f)$, where $Q_f=\{(x_1,x_2,\ldots,x_k,y)\mid f(x_1,\ldots,x_k)=y\}$.
We prove communication complexity lower bounds by reduction from different versions of the set disjointness problem. $\mathrm{DISJ}_{k,n}$ is a function $\{0,1\}^{kn}\to\{0,1\}$ such that for every $x_1,\ldots,x_k\in\{0,1\}^n$ the following holds:
\[\mathrm{DISJ}_{k,n}(x_1,\ldots,x_k)=\bigwedge_{j=1}^n\neg\left(\bigwedge_{i=1}^k(x_i)_j\right).\]
Let us define the communication promise problem $\mathrm{UDISJ}_{k,n}$ in the $k$-party NOF model. For each $i\in[k]$ the string $x_i$ is written on the forehead of the $i$th party, it is guaranteed that there exists at most one index $j\in[n]$ such that for every $i\in[k]$, $(x_i)_j=1$. The goal is to compute $\mathrm{DISJ}_{k,n}(x_1,\ldots,x_k)$.
\begin{authorclaim}{Theorem7 ([31,32])}
$R^{1/3}_{\mathrm{pub}}(\mathrm{UDISJ}_{k,n})=\Omega(\sqrt n/(2^kk))$.
\end{authorclaim}
For $k=2$ we omit the first index: $\mathrm{DISJ}_n=\mathrm{DISJ}_{2,n}$; in this case Theorem7 may be improved.
\begin{authorclaim}{Theorem8 ([16])}
$R^{1/3}_{\mathrm{pub}}(\mathrm{DISJ}_n)\ge R^{1/3}_{\mathrm{pub}}(\mathrm{UDISJ}_{2,n})=\Omega(n)$.
\end{authorclaim}
\paragraph{Proof complexity.} We consider refutational proof systems for the language of unsatisfiable CNF-formulas UNSAT. A refutation of $\phi\in\mathrm{UNSAT}$ in a proof system $\Pi$ is a sequence of Boolean functions (proof lines) such that each proof line either represents a clause of $\phi$ or derived from previous proof lines in the sequence via some sound inference rules. The last line of the proof is identically zero function. A proof system $\Pi$ is defined by a representation of proof lines and by a set of admissible inference rules. It is required that the inference rules are polynomially verifiable i.e. there exists an algorithm that checks whether it is legitimate to derive a line $L_0$ from the lines $L_1,\ldots,L_k$.
For example, in the Resolution proof lines are represented by clauses and the only inference rule is the resolution rule that allows deriving a clause $A\vee B$ from the clauses $A\vee x$ and $A\vee\neg x$.
The size of a proof is the total size of all representations of proof lines in the proof. The length of a proof is the number of proof lines in it.
A tree-like proof is such a proof that every its line can be used as a premise of a rule at most once. For each proof system, we can also consider its tree-like version where all proofs are constrained to be tree-like.
We also consider semantic refutational proof systems, where we drop the requirement for polynomial verification of inference rules i.e. we allow to derive any sound consequence from the premises. For such systems it is crucial to bound fan-in i.e. the number of the premises from which each proof line can be derived, otherwise, it would be possible to derive a contradiction from the clauses of the initial formula immediately. For example, it is well-known that Resolution is polynomially equivalent to a semantic proof system with fan-in 2 operating with clauses.
A lower bound on the proof size in a semantic proof system implies a lower bound on the proof size in its syntactic counterpart because a syntactic proof is always a semantic proof that operates with the same class of proof lines.
We define semantic $\mathrm{Res}(\oplus)$ as a semantic proof system with fan-in 2 that operates with linear clauses. A linear clause is a disjunction of linear equations over $\mathbb F_2$: $\bigvee_{i=1}^k(f_i=a_i)$, where $f_i$ is a linear form over $\mathbb F_2$ and $a_i\in\mathbb F_2$. Notice that an ordinary clause $\bigvee_{i\in P}x_i\vee\bigvee_{j\in N}\neg x_j$ can be represented by the linear clause $\bigvee_{i\in P}(x_i=1)\vee\bigvee_{j\in N}(x_j=0)$. For definition of syntactic version of $\mathrm{Res}(\oplus)$ we refer to [15]; it is also proved there that syntactic and semantic $\mathrm{Res}(\oplus)$ are polynomially equivalent.
We define semantic $\mathrm{Res}(\mathrm{PC}_d)$ as a semantic proof system with fan-in 2 that operates with disjunctions of equations of type $f=0$, where $f$ is a degree-$d$ polynomial over $\mathbb F_2$. Notice that semantic $\mathrm{Res}(\mathrm{PC}_1)$ is exactly semantic $\mathrm{Res}(\oplus)$. For the definition of the syntactic version of $\mathrm{Res}(\mathrm{PC}_d)$ we refer to [19].
Following [1] we define $\mathrm{Th}(k)$ as a semantic proof system with fan-in 2 that operates with polynomial inequalities $g\ge0$, where $g$ is a polynomial of degree at most $k$ with integer coefficients and Boolean variables. A clause $\bigvee_{i\in P}x_i\vee\bigvee_{j\in N}\neg x_j$ can be represented by an inequality $\sum_{i\in P}x_i+\sum_{j\in N}(1-x_j)-1\ge0$.
\paragraph{Proof complexity and communication complexity.} For an unsatisfiable CNF-formula $\phi$ we define the communication problem $\operatorname{Search}(\phi)$. Search$(\phi)$ is the following problem: given an assignment of the variables of the unsatisfiable CNF $\phi$, find a clause that is falsified by this assignment. It is assumed that variables of $\phi$ are somehow partitioned between the parties.
Following the paper [9] we consider a semantic proof system $\mathrm T^{\mathrm{cc}}(k,c)$ that models many interesting syntactic and semantic proof systems. The proof lines in $\mathrm T^{\mathrm{cc}}(k,c)$ can be arbitrary Boolean functions having the following property: for every proof line $C$ and every partition of variables of $C$ between $k$ parties, the NOF $k$-party randomized communication complexity of $C$ is at most $c$ w.r.t. this partition. We also define a semantic proof system $\mathrm T^{\mathrm{cc}}_{\mathrm{os}}(k,c)$ that is a subsystem of $\mathrm T^{\mathrm{cc}}(k,c)$ with the restriction that a communication protocol for proof lines must have a one-sided error: if the value of a proof line is zero, then the protocol should return zero with probability 1.
For example, $\mathrm T^{\mathrm{cc}}(2,2)$ simulates Resolution; $\mathrm T^{\mathrm{cc}}(2,O(1))$ simulates $\mathrm{Res}(\oplus)$ [22,15]; $\mathrm T^{\mathrm{cc}}(k,O(k^3\log^2n))$, where $n$ is the number of variables in a refuted formula, simulates $\mathrm{Th}(k-1)$ [9]. In Section3 we show that $\mathrm T^{\mathrm{cc}}_{\mathrm{os}}(d+1,O(1))$ simulates $\mathrm{Res}(\mathrm{PC}_d)$.
The following connection between the communication complexity of Search$(\phi)$ and tree-like proof complexity of $\phi$ is known.
\begin{authorclaim}{Lemma9 ([1,9])}
If a CNF formula $\phi$ has a tree-like $\mathrm T^{\mathrm{cc}}(k,c)$ refutation of length $\ell$ then, over any $k$-partition of the variables, there is a randomized bounded-error $k$-party NOF protocol for Search$(\phi)$ with communication cost $O(c\cdot\log\ell\log\log\ell)$.
\end{authorclaim}
In Section3 we show that for $\mathrm T^{\mathrm{cc}}_{\mathrm{os}}(k,c)$ the bound can be improved, see Remark14.
\paragraph{Basic formulas.} A CNF formula $\mathrm{PHP}_n^m$ encodes the pigeonhole principle; $\mathrm{PHP}_n^m$ states that it is possible to put $m$ pigeons into $n$ holes such that every pigeon flies to at least one hole and at most one pigeon flies to each hole. $\mathrm{PHP}_n^m$ depends on variables $p_{i,j}$ for $i\in[m]$ and $j\in[n]$ and $p_{i,j}=1$ iff the $i$-th pigeon flies to the $j$-th hole. $\mathrm{PHP}_n^m$ is the conjunction of $m(m-1)n/2$ hole axioms and $m$ pigeons axioms. For every $i\in[m]$ $\mathrm{PHP}_n^m$ contains a pigeon axiom $(p_{i,1}\vee p_{i,1}\vee\cdots\vee p_{i,n})$. And for every $j\in[n]$ and every $k\ne\ell\in[n]$, $\mathrm{PHP}_n^m$ contains a hole axiom $(\neg p_{k,j}\vee\neg p_{\ell,j})$. $\mathrm{PHP}_n^m$ is unsatisfiable iff $m>n$.
For an undirected graph $G(V,E)$, the formula $\mathrm{PM}_G$ encodes in CNF that $G$ has a perfect matching. The formula $\mathrm{PM}_G$ has $|E|$ variables, each of them corresponds to an edge of $G$, $x_e$ is the variable corresponding to $e\in E$.
\[\mathrm{PM}_G=\bigwedge_{v\in V}\left(\left(\bigvee_{e\text{ is incident to }v}x_e\right)\wedge\bigwedge_{e_1\ne e_2\text{ are incident to }v}(\neg x_{e_1}\vee\neg x_{e_2})\right).\]
$\mathrm{PM}_G$ is unsatisfiable iff $G$ does not have a perfect matching.
\begin{authorclaim}{Theorem10 ([26])}
Let $G$ be a graph with $n$ vertices, which has no perfect matching. Then the formula $\mathrm{PM}_G$ has a tree-like $\mathrm{Res}(\oplus)$ refutation of size $2^{O(n)}$.
\end{authorclaim}
\begin{authorclaim}{Proposition11 ([14])}
Let $G$ be a graph with an odd number of vertices. Then the formula $\mathrm{PM}_G$ has a tree-like $\mathrm{Res}(\oplus)$ refutation of size $\operatorname{poly}(n)$.
\end{authorclaim}
The binary pigeonhole principle $\mathrm{BPHP}_{2^\ell}^m$ states that there are $m$ different $\ell$-bit binary strings $s_1,s_2,\ldots,s_m$. $\mathrm{BPHP}_{2^\ell}^m$ has $m\ell$ variables corresponding to the bits of $s_i$ for $i\in[m]$. Then $\mathrm{BPHP}_{2^\ell}^m=\bigwedge_{i\ne j\in[m]}s_i\ne s_j$, where the predicate $s_i\ne s_j$ is encoded as a $2\ell$-CNF formula of size $2^\ell$ as follows: $\bigwedge_{\alpha\in\{0,1\}^\ell}(s_i\ne\alpha\vee s_j\ne\alpha)$; notice that the predicate $(s_i\ne\alpha\vee s_j\ne\alpha)$ can be represented by a clause with $2\ell$ literals. If $m>2^\ell$, then the formula $\mathrm{BPHP}_{2^\ell}^m$ is unsatisfiable.
Let $\phi$ be a CNF formula with $n$ variables, and $g:\{0,1\}^k\to\{0,1\}$ be a Boolean function. Then $\phi\circ g$ denotes a CNF formula on $kn$ variables that represents $\phi(g(\vec x_1),g(\vec x_2),\ldots,g(\vec x_n))$, where $\vec x_i$ denotes a vector of $k$ new variables. $\phi\circ g$ is constructed by applying the substitution to every clause $C$ of $\phi$ and converting the resulting function $C\circ g$ to CNF in some fixed way.

\paragraph{Original introduction: proof-complexity consequences and scope}
A lower bound on the tree-like communication complexity of $\operatorname{Search}(\phi)$ implies a lower bound on the size of tree-like refutations in the corresponding proof system, and on the size of dag-like refutations using small space (size-space tradeoffs [9,12]). The usual translation from tree-like refutations of size $S$, or small-space realizations of dag-like refutations of size $S$, gives communication cost $O(\log S\log\log S\cdot c)$, with $c$ the communication cost of a proof line; sometimes it improves to $O(\log S\cdot c)$.
Randomized two-party protocols correspond to $\mathrm{Res}(\oplus)$ [15], whose lines are disjunctions of linear equations over $\mathbb F_2$. Real $k$-party NOF protocols correspond to semantic $\mathrm{Th}(k-1)$ [1], whose lines are integer-coefficient polynomial inequalities of degree at most $k-1$ on Boolean variables. Randomized $k$-party NOF protocols correspond to semantic $\mathrm T^{\mathrm{cc}}(k,c)$, whose lines are arbitrary predicates of randomized NOF communication cost at most $c$. For small $c$ this simulates $\mathrm{Th}(k-1)$ and $\mathrm{Res}(\mathrm{PC}_{k-1})$; the latter's lines are disjunctions of polynomial equations over $\mathbb F_2$ of bounded degree.
Theorem4 yields the lower bound $\exp(2^{\Omega(n/k)})$ on the size of tree-like $\mathrm T^{\mathrm{cc}}(k,c)$ (and $\mathrm{Th}(k-1)$) refutations of $\mathrm{BPHP}_{2^n}^{2^n+2^{n+k-\lfloor n/k\rfloor}}$.
Hrubes and Pudlak [10] obtained lower bounds for dag-like two-party real protocols with the same variable partition and arbitrary $m>2^\ell$; the results are formally incomparable because the present randomized model is stronger and multiparty, while [10] covers dag-like protocols and arbitrary weak binary PHP.
The randomized communication lower bound is nontrivial for $k\le\log^{1-\varepsilon}n$ with $\varepsilon>0$, yielding superpolynomial tree-like $\mathrm{Th}(k)$ lower bounds for $\mathrm{BPHP}_n^m$. Proposition6 gives short tree-like $\mathrm{Th}(\log n)$ refutations. Together with [10] it implies tree-like $\mathrm{Th}(\log n)$ cannot be simulated by dag-like $\mathrm{Th}(1)$.
Theorem4 and Proposition6 imply that $\mathrm{BPHP}_{2^\ell}^{2^\ell+1}$ separates tree-like $\mathrm{Th}(\log n)$ from tree-like $\mathrm{Th}(k)$ for $k\le\log^{1-\varepsilon}n$. Here $n=(2^\ell+1)\ell$ is the number of variables; Proposition6 gives polynomial size, while Theorem4 gives at least $\exp(\exp(\Omega(\log^\varepsilon n)))$, which is superpolynomial.
The abstract states rank lower bounds among the general consequences of randomized NOF search bounds; these are contextual consequences rather than a separately counted claim.

% Faithful formal-publication transcription: complete Section3.
\section*{3 Communication protocols from tree-like Res(PCd) proofs}
Let $\phi$ be an unsatisfiable CNF formula with $n$ variables. Let us define the communication problem $\oplus_k\operatorname{Search}(\phi)$ with $k$ parties as follows. Assume that the $i$th party has an assignment $\alpha_i\in\{0,1\}^n$ written on the forehead. They aim to find a clause of $\phi$ falsified by the assignment $\sum_{i=1}^k\alpha_i$ (all sums of boolean vectors are computed modulo 2).
\begin{authorclaim}{Lemma1}
Let $\phi$ be an unsatisfiable CNF formula. If there exists a tree-like $\mathrm{Res}(\mathrm{PC}_d)$ proof of $\phi$ of length $m$, then $R^{1/3}_{\mathrm{pub}}(\oplus_{d+1}\operatorname{Search}(\phi))=O(d\cdot\log m)$.
\end{authorclaim}
A slightly weaker version of the following lemma was implicitly proved in [15]:
\begin{authorclaim}{Lemma12 (see proof of Theorem3.11 from [15])}
Let $T$ be a binary tree with $m$ vertices such that the $i$th vertex is labeled with $a_i\in\{0,1\}$ with the hereditary property: for each inner vertex $i$ with direct descendants $c_1$ and $c_2$, if $a_i=1$, then $a_{c_1}=1$ or $a_{c_2}=1$. We also assume that if $r$ is the root of $T$, then $a_r=1$. Assume that we have a one-sided bounded error oracle access to $a_i$ i.e. if we request a value of $a_i$ and $a_i=0$ we get 1 with probability at most $1/2$ and0 with probability at least $1/2$; if $a_i=1$ we get 1 with probability 1. Then there exists an algorithm $A$ that with probability at least $2/3$ returns a leaf $\ell$ of $T$ with $a_\ell=1$ and makes $O(\log m)$ oracle queries to $a_1,\ldots,a_m$.
\end{authorclaim}
Proof. See AppendixA.
\begin{proof}[Proof of Lemma1]
Let $F_1,\ldots,F_m$ be a tree-like $\mathrm{Res}(\mathrm{PC}_d)$-refutation of $\phi$ with the underlying tree $T$, where vertices of $T$ are identified with $[m]$. Then the leaves of $T$ correspond to the clauses of $\phi$ and $m$ is the root of $T$.
Let $\alpha_1,\ldots,\alpha_{d+1}$ be the assignments written on the foreheads of $d+1$ parties. Let $\alpha=\sum_{i=1}^{d+1}\alpha_i$. Let $a_i=1$ iff $\alpha$ falsifies $F_i$ for $i\in[m]$. Then $a_m=1$ since $F_m$ is identically false. For any inner node $v$ of $T$, if $a_v=1$ then for the direct descendants of $v$, $c_1$ and $c_2$ either $a_{c_1}=1$ or $a_{c_2}=1$. In the next paragraphs we show that for any $i\in[m]$ there exists a NOF $(d+1)$-party protocol that computes $a_i$ given that for each $j\in[d+1]$ the $j$th party has $\alpha_j$ written on their forehead such that
\begin{itemize}
\item the protocol transmits $O(d)$ bits;
\item the protocol has one-sided bounded error: if $a_i=1$ then the protocol returns 1 with probability 1 and if $a_i=0$ the protocol returns 0 with probability at least $1/2$.
\end{itemize}
Then we use this protocol to compute $a_i$ as an oracle in the algorithm given by Lemma12 and thus show that there is a NOF $(d+1)$-party protocol computing $\oplus_{d+1}\operatorname{Search}(\phi)$ with communication cost $O(d\log m)$.
Now we show that for every $\ell\in[m]$, $F_\ell(\alpha)$ can be computed by a $(d+1)$-party NOF protocol with one-sided error using $O(d)$ bits of communication. Let $F_\ell=\bigvee_{j=1}^t(f_j=1)$, where $f_1,\ldots,f_t$ are polynomials over $\mathbb F_2$ of degree at most $d$. Let $z_1,\ldots,z_n$ be the variables of $\phi$. Let us introduce new variables $y_{1,1},\ldots,y_{1,n},\ldots,y_{d+1,1},\ldots,y_{d+1,n}$ and assume that for each $i\in[d+1]$ the $i$th party has the value of variables $y_{i,1},y_{i,2},\ldots,y_{i,n}$ written on the forehead or in other words $\alpha_i$ assigns values of $y_{i,1},y_{i,2},\ldots,y_{i,n}$. Let $\bar f_j$ denote $f_j$ after substitution $z_\ell:=y_{1,\ell}+y_{2,\ell}+\cdots+y_{d+1,\ell}$ for $\ell\in[n]$; $j\in[t]$. Since for all $j\in[t]$, $\deg\bar f_j=\deg f_j\le d$, we can represent $\bar f_j=\bar f_j^{(1)}+\cdots+\bar f_j^{(d+1)}$ such that $\bar f_j^{(s)}$ does not contain variables $y_{s,1},\ldots,y_{s,n}$ for each $s\in[d+1]$. Then the $i$th party can compute $\bar f_1^{(i)}(\alpha_1,\ldots,\alpha_{d+1}),\ldots,\bar f_t^{(i)}(\alpha_1,\ldots,\alpha_{d+1})$. Notice that $F_\ell=\neg(\bigwedge_{j=1}^t(f_j=0))$.
The final step of the protocol exploits the idea used to construct a short randomized communication protocol for equality. Take a random uniformly distributed vector $(e_1,\ldots,e_t)\in\mathbb F_2^t$. Then all parties compute
\[\sum_{j=1}^te_jf_j(\alpha)=\sum_{i=1}^{d+1}\underbrace{\sum_{j=1}^te_j\bar f_j^{(i)}}_{i\text{th party}}\]
with $O(d)$ bits of communication and the protocol halts. To bound the error probability we use the following well-known statement:
\begin{authorclaim}{Proposition13 (Random subsum principle)}
For any $x\in\mathbb F_2^k\setminus\{0^k\}$,
\[\Pr_{y\leftarrow U(\mathbb F_2^k)}\left[\sum_{i=1}^ky_ix_i=1\right]=\frac12.\]
\end{authorclaim}
If $F_\ell(\alpha)=1$ then $\Pr[\sum_{j=1}^te_jf_j(\alpha)\ne0]=1/2$ by the random subsum principle. If $F_\ell(\alpha)=0$, then $\Pr[\sum_{j=1}^te_jf_j(\alpha)=0]=1$.
\end{proof}
\begin{authorclaim}{Remark14}
Similarly to the proof of Lemma1 one can prove that if an unsatisfiable CNF formula $\phi$ has a tree-like $\mathrm T^{\mathrm{cc}}_{\mathrm{os}}(k,c)$ refutation of length $\ell$, then for any $k$-partition of the variables, there is a randomized bounded-error $k$-party NOF protocol for Search$(\phi)$ with communication cost $O(c\log\ell)$. Thus, the bound from Lemma9 can be slightly improved in the case of one-sided error.
\end{authorclaim}

\begin{authorclaim}{Theorem 3}
Let $\ell$ and $k$ be natural numbers such that $2\le k\le\ell-7$. Then
\[
R^{1/3}_{\mathrm{pub}}(\oplus_k\operatorname{Search}(\mathrm{BPHP}_{2^\ell}^{2^\ell+2^k}))
=\Omega\left(\frac{2^{\ell/2}}{k2^{3k/2}}\right).
\]
For $k=2$ the stronger bound holds:
$R^{1/3}_{\mathrm{pub}}(\oplus_2\operatorname{Search}(\mathrm{BPHP}_{2^\ell}^{2^\ell+4}))=\Omega(2^\ell)$.
\end{authorclaim}

We consider a combinatorial analogue of the communication problem $\oplus_k\operatorname{Search}(\mathrm{BPHP}_{2^\ell}^m)$. Assume that each of $k$ parties gets $m$ binary strings from $\{0,1\}^\ell$, where $m>2^\ell$. The $i$th party has numbers $a_{i,1},\ldots,a_{i,m}\in\{0,1\}^\ell$ on their forehead. Based on these strings we form $m$ vectors from $\mathbb F_2^\ell$: $x_1,\ldots,x_m$, where $x_j=\sum_{i=1}^ka_{i,j}$. The goal is to find different $t,s\in[m]$ such that $x_t=x_s$. We denote this problem by $\oplus_k\mathrm{BPHP}_{2^\ell}^m$. It is straightforward that $R^{1/3}_{\mathrm{pub}}(\oplus_k\operatorname{Search}(\mathrm{BPHP}_{2^\ell}^m))\ge R^{1/3}_{\mathrm{pub}}(\oplus_k\mathrm{BPHP}_{2^\ell}^m)$, hence it is sufficient to prove a lower bound on the latter.
\begin{authorclaim}{Theorem 18}
Let $\ell,k$ be natural numbers such that $2\le k\le\ell-7$. Then
\[
R^{1/3}_{\mathrm{pub}}(\oplus_k\mathrm{BPHP}_{2^\ell}^{2^\ell+2^k})
=\Omega\big(R^{1/3}_{\mathrm{pub}}(\mathrm{UDISJ}_{k,2^{\ell-k}-1})-\ell\big).
\]
\end{authorclaim}
\begin{authorclaim}{Corollary 19}
$R^{1/3}_{\mathrm{pub}}(\oplus_k\mathrm{BPHP}_{2^\ell}^{2^\ell+2^k})=\Omega(2^{\ell/2}/(k2^{3k/2}))$. For $k=2$: $R^{1/3}_{\mathrm{pub}}(\oplus_2\mathrm{BPHP}_{2^\ell}^{2^\ell+4})=\Omega(2^\ell)$.
\end{authorclaim}
\begin{proof}[Proof of Corollary 19]
Follows from Theorem 18 and Theorem 7; for $k=2$ apply Theorem 8.
\end{proof}
Theorem 3 immediately follows from Corollary 19.
\subsection*{5.1 Warm-up example}
Consider $\mathrm{Distinct}_{k,\ell}$: each of $k$ parties has a matrix from $\mathbb F_2^{2^\ell\times\ell}$ on their forehead. The goal is to determine whether all rows of the sum of these matrices are distinct. A version without the xor-gadget is referred to as Element Distinctness (ED) [25].
\begin{authorclaim}{Proposition 20}
$R^{1/3}_{\mathrm{pub}}(\mathrm{Distinct}_{k,\ell})\ge R^{1/3}_{\mathrm{pub}}(\mathrm{UDISJ}_{k,2^{\ell-k}})$.
\end{authorclaim}
Let $S_k$ denote the set of matrices from $\{0,1\}^{2^k\times k}$ with all distinct rows. Let $K_k\in\{0,1\}^{2^k\times k}$ have $i$th row $\operatorname{bin}_k(i-1-((i-1)\bmod2))$, i.e. rows $\operatorname{bin}_k(0),\operatorname{bin}_k(0),\operatorname{bin}_k(2),\operatorname{bin}_k(2),\ldots,\operatorname{bin}_k(2^{k-1}-2),\operatorname{bin}_k(2^{k-1}-2)$. Every row starts with zero and appears exactly twice.
The following combinatorial lemma is proved in Subsection 5.4.
\begin{authorclaim}{Lemma 21}
There exist matrices $A_1(0),A_1(1),\ldots,A_k(0),A_k(1)\in\mathbb F_2^{2^k\times k}$ such that $\sum_{i=1}^kA_i(1)=K_k$ and for all $b_1,\ldots,b_k\in\{0,1\}$, if $\bigwedge_{i=1}^kb_i=0$, then $\sum_{i=1}^kA_i(b_i)\in S_k$.
\end{authorclaim}
\begin{proof}[Proof of Proposition 20]
Let $(x_{i,1},\ldots,x_{i,2^{\ell-k}})$ be an input of the $i$th party of the problem $\mathrm{UDISJ}_{k,2^{\ell-k}}$. For all $i\in[k]$ we construct a matrix $D_i$ of size $2^\ell\times\ell$ and put it on the forehead of the $i$th party. Let $A_i(b)$ for $i\in[k]$, $b\in\{0,1\}$ be matrices of size $2^k\times k$ from Lemma21. Let $J_t$ for $t\in[1,\ldots,2^{\ell-k}]$ be a matrix of size $2^k\times(\ell-k)$ such that all its rows are equal to $\operatorname{bin}_{\ell-k}(t-1)$. Let us define
\[
D_1:=\begin{pmatrix}J_1&A_1(x_{1,1})\\\vdots&\vdots\\J_j&A_1(x_{1,j})\\\vdots&\vdots\\J_{2^{\ell-k}}&A_1(x_{1,2^{\ell-k}})\end{pmatrix};
\quad
D_i:=\begin{pmatrix}0_{2^k\times(\ell-k)}&A_i(x_{i,1})\\\vdots&\vdots\\0_{2^k\times(\ell-k)}&A_i(x_{i,j})\\\vdots&\vdots\\0_{2^k\times(\ell-k)}&A_i(x_{i,2^{\ell-k}})\end{pmatrix}
\quad(i\in\{2,\ldots,k\}).
\]
By Lemma21, the matrix $D_1+D_2+\cdots+D_k$ has the following property: for all $j\in[2^{\ell-k}]$, its submatrix formed by the rows with numbers from $[2^k\cdot(j-1)+1,2^k\cdot j]$ has two equal rows if and only if $x_{1,j}=x_{2,j}=\cdots=x_{k,j}=1$. Thus, the communication complexity of $\mathrm{UDISJ}_{k,2^{\ell-k}}$ is at most the communication complexity of $\mathrm{Distinct}_{k,\ell}$.
\end{proof}
\subsection*{5.2 Proof of Theorem 18}
We modify Proposition20 to reduce $\mathrm{UDISJ}_{k,2^{\ell-k}-1}$ to $\oplus_k\mathrm{BPHP}_{2^\ell}^{2^\ell+2^k}$ by adding fake rows not corresponding to the disjointness input and randomizing to hide them among other rows.
\begin{proof}[Proof of Theorem 18]
For $N>2^\ell$, define $\mathrm{ROW}\oplus_k\mathrm{BPHP}_{2^\ell}^N$: party $i$ has $M_i\in\mathbb F_2^{N\times\ell}$ on their forehead and the goal is to find the value of a row appearing at least twice in $M_1+\cdots+M_k$, rather than the indices of equal rows.
\begin{authorclaim}{Claim 22}
If $R^{1/3}(\oplus_k\mathrm{BPHP}_{2^\ell}^N)\le t$, there exists a protocol $P$ for $\mathrm{ROW}\oplus_k\mathrm{BPHP}_{2^\ell}^N$ using $O(t+\ell)$ bits that returns a correct answer or $\bot$, with $\Pr[P(M_1,\ldots,M_k)=\bot]\le1/100$ for every input.
\end{authorclaim}
\begin{proof}
$P$ executes a randomized protocol for $\oplus_k\mathrm{BPHP}_{2^\ell}^N$ and verifies its answer by transferring additional $O(\ell)$ bits. The probability of failure can be reduced by repetition.
\end{proof}
We describe a protocol for $\mathrm{UDISJ}_{k,2^{\ell-k}-1}$ using $P$ for $\mathrm{ROW}\oplus_k\mathrm{BPHP}_{2^\ell}^{2^\ell+2^k}$. Inputs are $x_1,\ldots,x_k\in\{0,1\}^{2^{\ell-k}-1}$; $x_{i,j}$ is bit $j$ of $x_i$ and $\vec x=(x_1,\ldots,x_k)$.
\paragraph{Important matrices}
Let $\gamma:[2^{\ell-k}-1]\cup\{*\}\to\{0,1\}^{\ell-k}$ be bijective. Define matrices $D_1(x_1,\gamma),D_2(x_2),\ldots,D_k(x_k)$ of size $(2^\ell+2^k)\times\ell$. Let $A_i(b)$ be from Lemma21; for every $t\in\{0,1\}^{\ell-k}$ let $J_t$ have size $2^k\times(\ell-k)$ with every row $t$. Fix $W\in S_k$. Define
\[
D_1(x_1,\gamma):=\begin{pmatrix}
J_{\gamma(1)}&A_1(x_{1,1})\\\vdots&\vdots\\J_{\gamma(j)}&A_1(x_{1,j})\\\vdots&\vdots\\J_{\gamma(2^{\ell-k}-1)}&A_1(x_{1,2^{\ell-k}-1})\\J_{\gamma(*)}&W\\J_{\gamma(*)}&W
\end{pmatrix};
\]
for $i\in[k]\setminus\{1\}$,
\[
D_i(x_i):=\begin{pmatrix}
0_{2^k\times(\ell-k)}&A_i(x_{i,1})\\\vdots&\vdots\\0_{2^k\times(\ell-k)}&A_i(x_{i,j})\\\vdots&\vdots\\0_{2^k\times(\ell-k)}&A_i(x_{i,2^{\ell-k}-1})\\0_{2^k\times(\ell-k)}&0_{2^k\times k}\\0_{2^k\times(\ell-k)}&0_{2^k\times k}
\end{pmatrix}.
\]
The last $2^{k+1}$ rows of $D_1(x_1,\gamma)$ contain each row exactly twice.
Define $H_{\vec x}(\gamma):=D_1(x_1,\gamma)+D_2(x_2)+\cdots+D(x_k)$. Lemma21 gives the following property in the standard basis.
\begin{authorclaim}{Definition 23}
Let $M\in\mathbb F_2^{(2^k+2^\ell)\times\ell}$, $\gamma:[2^{\ell-k}-1]\cup\{*\}\to\{0,1\}^{\ell-k}$ bijective, and $e_1,\ldots,e_\ell$ a basis in $\mathbb F^\ell$. $M$ satisfies the key property w.r.t. $(\gamma,\vec x)$ in this basis if:
\begin{itemize}
\item For a row $s$ among the last $2^{k+1}$ rows, the first $\ell-k$ basis coordinates are $\gamma(*)_1,\ldots,\gamma(*)_{\ell-k}$ and $s$ appears exactly twice.
\item For a row $s$ with index in $[2^k(i-1)+1;2^ki]$, $i\in[2^{\ell-k}-1]$, the first $\ell-k$ basis coordinates are $\gamma(i)_1,\ldots,\gamma(i)_{\ell-k}$. If $\bigwedge_{j=1}^kx_{i,j}=0$ it appears exactly once; if $\bigwedge_{j=1}^kx_{i,j}=1$ it appears exactly twice and its $(\ell-k+1)$th basis coordinate is 0.
\end{itemize}
\end{authorclaim}
Consider an invertible matrix $E\in\mathbb F_2^{\ell\times\ell}$. Let $e_1,e_2,\ldots,e_\ell$ be the rows of $E$. Since $E$ is invertible, $e_1,e_2,\ldots,e_\ell$ form a basis. Let us define $C_{\vec x}(\gamma,E):=H(\vec x,\gamma)E$. Rows of $C_{\vec x}(\gamma,E)$ can be viewed as vectors with coordinates in the basis $e_1,e_2,\ldots,e_\ell$ corresponding to the rows of $H(\vec x,\gamma)$. Hence, $C_{\vec x}(\gamma,E)$ satisfies the key property w.r.t. $(\gamma,\vec x)$ in the basis $(e_1,e_2,\ldots,e_\ell)$.
Define $\operatorname{Fake}(\gamma,E)\subseteq\mathbb F_2^\ell$ to be the set of the last $2^{k+1}$ rows of $C_{\vec x}(\gamma,E)$. This set does not depend on $\vec x$. Its rows appear exactly twice.
\paragraph{Random variables}
Random variables are bold. Public randomness consists of:
\begin{itemize}
\item Uniform random bijection $\boldsymbol\gamma:[2^{\ell-k}-1]\cup\{*\}\to\{0,1\}^{\ell-k}$.
\item Uniform random invertible matrix $\mathbf E\in\mathbb F_2^{\ell\times\ell}$.
\item Uniform random permutation $\boldsymbol\pi$ of $[2^\ell+2^k]$, with corresponding permutation matrix $M_{\boldsymbol\pi}$: $(M_{\boldsymbol\pi})_{i,j}=1\Longleftrightarrow\boldsymbol\pi(i)=j$.
\item Uniform random matrices $\boldsymbol\Delta_1,\ldots,\boldsymbol\Delta_k\in\mathbb F_2^{(2^\ell+2^k)\times\ell}$ subject to $\sum_i\boldsymbol\Delta_i=0$.
\end{itemize}
Define $\mathbf P_i=M_{\boldsymbol\pi}D_i(x_i)\mathbf E+\boldsymbol\Delta_i$ for $i\ge2$ and $\mathbf P_1=M_{\boldsymbol\pi}D_1(x_1,\boldsymbol\gamma)\mathbf E+\boldsymbol\Delta_1$. The addition of $\boldsymbol\Delta_i$ makes each $\mathbf P_i$ indistinguishable from a random matrix. Moreover, $\sum_i\mathbf P_i=M_{\boldsymbol\pi}C_{\vec x}(\boldsymbol\gamma,\mathbf E)$, obtained by permuting its rows. Let $N$ be a constant to be chosen later.
\begin{authorclaim}{Algorithm 1: Protocol T solving $\mathrm{UDISJ}_{k,2^{\ell-k}-1}$}
Input $x_1,\ldots,x_k\in\{0,1\}^{2^{\ell-k}-1}$; $x_i$ is on party $i$'s forehead.
\begin{algorithmic}
\State $Z:=\varnothing$
\Loop\Statex repeat $N$ times
\State Sample $\pi\leftarrow\boldsymbol\pi$, $E\leftarrow\mathbf E$, $\gamma\leftarrow\boldsymbol\gamma$, $\vec\Delta\leftarrow\boldsymbol{\vec\Delta}$ \Comment{Fresh public bits}
\State $P_1:=M_\pi D_1(x_1,\gamma)E+\Delta_1$
\Statex $\triangleright$ Computable by parties $2,\ldots,k$.
\State $P_i:=M_\pi D_i(x_i)E+\Delta_i$ for $i\ge2$
\Statex $\triangleright$ Computable by all parties except $i$.
\State $z:=P(P_1,\ldots,P_k)$
\Statex $\triangleright$ Fresh random bits for $P$; $P_i$ is on party $i$'s forehead.
\State $Z:=Z\cup\{z\}$
\EndLoop
\If{$Z=\{\bot\}$}\State \Return $\bot$
\ElsIf{$Z\setminus\{\bot\}\subseteq\operatorname{Fake}(\gamma,E)$}\State \Return 1
\Statex $\triangleright$ Most likely there are no repeated rows except Fake; the key property gives $\mathrm{DISJ}(x_1,\ldots,x_k)=1$.
\EndIf
\State \Return 0
\end{algorithmic}
\end{authorclaim}
\paragraph{Protocol analysis}
Let us analyze the protocol $T$. Since it executes the protocol $P$ a constant number of times, $T$ transmits $O(t+\ell)$ bits. Assume that $x_1,x_2,\ldots,x_k$ is a 1-instance of $\mathrm{UDISJ}_{k,2^\ell+2^k}$. Then by the key property of $C_{\vec x}(\boldsymbol\gamma,\mathbf E)$ all repeated rows of $\sum_{i=1}^k\mathbf P_i$ are in $\operatorname{Fake}(\boldsymbol\gamma,\mathbf E)$, hence the protocol $T$ returns either $\bot$ or the correct answer. Since $P$ is executed $N$ times independently, the probability that $Z=\{\bot\}$ is at most $1/100^N$, hence $T$ returns1 with probability at least $1-1/100^N$.
The rest of the proof is devoted to the analysis of the case, where $x_1,x_2,\ldots,x_k$ is a 0-instance of $\mathrm{UDISJ}_{k,2^\ell+2^k}$. This is the most technically involved part of the proof. So it is a good point to give a large scale overview of the further proof strategy. Our goal is to show that if $x_1,x_2,\ldots,x_k$ is a 0-instance of $\mathrm{UDISJ}_{k,2^\ell+2^k}$, then the probability that $P(\mathbf P_1,\ldots,\mathbf P_k)$ returns a value from $\operatorname{Fake}(\boldsymbol\gamma,\mathbf E)$ is bounded by some constant less than1.
The random variable $P(\mathbf P_1,\ldots,\mathbf P_k)$ depends on random bits used by the protocol $P$ and on random bits needed for sampling $\mathbf P_1,\ldots,\mathbf P_k$. Let $R$ denote the set of all random strings used by the protocol $P$ (i.e. we assume that $P$ sample a random string from $R$ and use it as public randomness) and $S$ denote the set of all random strings used for sampling $\mathbf P_1,\ldots,\mathbf P_k$. We would like to construct two bijections $\alpha$ and $\beta$ on the set $S$ such that for every $s\in S$ the following two properties hold.
\begin{enumerate}
\item The three values of random variable $(\mathbf P_1,\ldots,\mathbf P_k)$ sampled using three strings $s,\alpha(s)$ and $\beta(s)$ as a random source, are the same.
\item Let $(\gamma,E)$, $(\gamma_\alpha,E_\alpha)$ and $(\gamma_\beta,E_\beta)$ be values of the random variable $(\boldsymbol\gamma,\mathbf E)$ that is sampled using three strings $s,\alpha(s)$ and $\beta(s)$ as a random source. Then $\operatorname{Fake}(\gamma,E)\cap\operatorname{Fake}(\gamma_\alpha,E_\alpha)\cap\operatorname{Fake}(\gamma_\beta,E_\beta)=\varnothing$.
\end{enumerate}
Consider arbitrary strings $r\in R$ and $s\in S$. The first property implies that for random variables sampled using strings $(r,s)$, $(r,\alpha(s))$ and $(r,\beta(s))$ as a random source values of $P(\mathbf P_1,\ldots,\mathbf P_k)$ are the same. The second property implies that for at least one of this cases this value does not belong to $\operatorname{Fake}(\boldsymbol\gamma,\mathbf E)$. Then, using that $\alpha$ and $\beta$ are bijections, we get $\Pr[P(\mathbf P_1,\ldots,\mathbf P_k)\in\operatorname{Fake}(\boldsymbol\gamma,\mathbf E)]\le2/3$.
Since we have many random variables, it is a tedious task to construct such $\alpha$ and $\beta$. In order to simplify this task we slightly relax the properties. We will define bijections $\alpha$ and $\beta$ not on all strings $S$ but only on the part of bits corresponding to sampling of $\boldsymbol\gamma$ and $\mathbf E$. More precisely we will define two bijections $\alpha$ and $\beta$ on the set of values of the random variable $(\boldsymbol\gamma,\mathbf E)$. We relax the first property as follows:
\begin{enumerate}
\item[1'.] For every $\gamma$ and $E$ the three conditional distributions of the random variable $(\mathbf P_1,\ldots,\mathbf P_k)$ under the following three conditions coincide:
\begin{enumerate}\item $(\boldsymbol\gamma,\mathbf E)=(\gamma,E)$,\item $(\boldsymbol\gamma,\mathbf E)=\alpha(\gamma,E)$ and\item $(\boldsymbol\gamma,\mathbf E)=\beta(\gamma,E)$.\end{enumerate}
\end{enumerate}
Unfortunately, we were not able to construct such bijections on the set of all pairs $(\gamma,E)$. Thus we take a set $\Xi$ consisting $1-\delta$ fraction of all values of $(\boldsymbol\gamma,\mathbf E)$ and we will claim that $\alpha$ and $\beta$ are bijections on $\Xi$. Such relaxation will weaken the bound of the probability up to $2/3+\delta$. We formalize the requirements to $\Xi,\alpha$ and $\beta$ in Definition24. Then we verify in Claim25 that these requirements are sufficient to bound $\Pr[P(\mathbf P_1,\ldots,\mathbf P_k)\in\operatorname{Fake}(\boldsymbol\gamma,\mathbf E)]$. The construction of $\Xi,\alpha$ and $\beta$ is given in Subsection5.3.
\begin{authorclaim}{Definition 24}
For a 0-instance $\vec x$ of $\mathrm{UDISJ}_{k,2^{\ell-k}-1}$ and constant $1>\delta\ge0$, let $\Xi$ consist of pairs $(\gamma,E)$ with $\gamma:[2^{\ell-k}-1]\cup\{*\}\to\{0,1\}^{\ell-k}$ bijective and $E\in\mathbb F_2^{\ell\times\ell}$ invertible. Bijections $\alpha,\beta:\Xi\to\Xi$ form a $(1-\delta)$-symmetry randomness space if:
\begin{itemize}
\item (Largeness) $\Pr[(\boldsymbol\gamma,\mathbf E)\in\Xi]\ge1-\delta$.
\item (Difference) $\operatorname{Fake}(\gamma,E)\cap\operatorname{Fake}(\alpha(\gamma,E))\cap\operatorname{Fake}(\beta(\gamma,E))=\varnothing$ for every $(\gamma,E)\in\Xi$.
\item (Symmetry) $C_{\vec x}(\gamma,E)$, $C_{\vec x}(\alpha(\gamma,E))$, $C_{\vec x}(\beta(\gamma,E))$ differ only by row permutation.
\end{itemize}
\end{authorclaim}
\begin{authorclaim}{Claim 25}
Under Definition24, $\Pr[P(\mathbf P_1,\ldots,\mathbf P_k)\in\operatorname{Fake}(\boldsymbol\gamma,\mathbf E)]\le2/3+\delta$.
\end{authorclaim}
\begin{proof}
Write $\boldsymbol{\vec P}=(\mathbf P_1,\ldots,\mathbf P_k)$, $\boldsymbol{\vec\Delta}=(\boldsymbol\Delta_1,\ldots,\boldsymbol\Delta_k)$, $\vec D(\vec x,\gamma)=(D_1(x_1,\gamma),D_2(x_2),\ldots,D_k(x_k))$, so $\boldsymbol{\vec P}=\boldsymbol{\vec\Delta}+M_{\boldsymbol\pi}(\vec D(\vec x,\boldsymbol\gamma)\mathbf E)$.
Let $p=\Pr[P(\boldsymbol{\vec P})\in\operatorname{Fake}(\boldsymbol\gamma,\mathbf E)]$. Then
\begin{align*}
p&=\sum_{\gamma,E}\Pr[P(\boldsymbol{\vec\Delta}+M_{\boldsymbol\pi}(\vec D(\vec x,\gamma)E))\in\operatorname{Fake}(\gamma,E)]\Pr[\boldsymbol\gamma=\gamma,\mathbf E=E]\\
&\le\sum_{(\gamma,E)\in\Xi}\Pr[P(\boldsymbol{\vec\Delta}+M_{\boldsymbol\pi}(\vec D(\vec x,\gamma)E))\in\operatorname{Fake}(\gamma,E)]\Pr[\boldsymbol\gamma=\gamma,\mathbf E=E]+\delta .
\end{align*}
Notice that for fixed $\gamma,E$ the random variable $\boldsymbol{\vec\Delta}+M_{\boldsymbol\pi}\cdot(\vec D(\vec x,\gamma)\cdot E)$ is distributed uniformly on the set of tuples $(L_1,\ldots,L_k)$ of $k$ matrices from $\mathbb F_2^{(2^\ell+2^k)\times\ell}$ such that $\sum_{i=1}^kL_i$ differs from $C_{\vec x}(\gamma,E)$ only by a permutation of rows. Let $(\gamma_{\alpha^{-1}},E_{\alpha^{-1}})=\alpha^{-1}(\gamma,E)$. By the symmetry condition, matrices $C_{\vec x}(\gamma,E)$ and $C_{\vec x}(\gamma_{\alpha^{-1}},E_{\alpha^{-1}})$ differ only by permutation of rows. Thus, for every set $A$ the probability
\[
\Pr[P(\boldsymbol{\vec\Delta}+M_{\boldsymbol\pi}(\vec D(\vec x,\gamma)E))\in A]
=\Pr[P(\boldsymbol{\vec\Delta}+M_{\boldsymbol\pi}(\vec D(\vec x,\gamma_{\alpha^{-1}})E_{\alpha^{-1}}))\in A].
\]
Therefore
\begin{align*}
p&\le\sum_{(\gamma,E)\in\Xi}\Pr[P(\boldsymbol{\vec\Delta}+M_{\boldsymbol\pi}(\vec D(\vec x,\gamma_{\alpha^{-1}})E_{\alpha^{-1}}))\in\operatorname{Fake}(\gamma,E)]\Pr[\boldsymbol\gamma=\gamma,\mathbf E=E]+\delta\\
&=\sum_{(\gamma,E)\in\Xi}\Pr[P(\boldsymbol{\vec\Delta}+M_{\boldsymbol\pi}(\vec D(\vec x,\gamma)E))\in\operatorname{Fake}(\alpha(\gamma,E))]\Pr[(\boldsymbol\gamma,\mathbf E)=\alpha(\gamma,E)]+\delta\\
&=\sum_{(\gamma,E)\in\Xi}\Pr[P(\boldsymbol{\vec\Delta}+M_{\boldsymbol\pi}(\vec D(\vec x,\gamma)E))\in\operatorname{Fake}(\alpha(\gamma,E))]\Pr[(\boldsymbol\gamma,\mathbf E)=(\gamma,E)]+\delta\\
&=\Pr[P(\boldsymbol{\vec P})\in\operatorname{Fake}(\alpha(\boldsymbol\gamma,\mathbf E)),(\boldsymbol\gamma,\mathbf E)\in\Xi]+\delta .
\end{align*}
Analogously, $p\le\Pr[P(\boldsymbol{\vec P})\in\operatorname{Fake}(\beta(\boldsymbol\gamma,\mathbf E)),(\boldsymbol\gamma,\mathbf E)\in\Xi]+\delta$. Also the inequality $p\le\Pr[P(\boldsymbol{\vec P})\in\operatorname{Fake}(\boldsymbol\gamma,\mathbf E),(\boldsymbol\gamma,\mathbf E)\in\Xi]+\delta$ follows by the largeness condition. Then,
\begin{align*}
3(1-p)&\ge\Pr[P(\boldsymbol{\vec P})\notin\operatorname{Fake}(\beta(\boldsymbol\gamma,\mathbf E))\vee(\boldsymbol\gamma,\mathbf E)\notin\Xi]\\
&\quad+\Pr[P(\boldsymbol{\vec P})\notin\operatorname{Fake}(\alpha(\boldsymbol\gamma,\mathbf E))\vee(\boldsymbol\gamma,\mathbf E)\notin\Xi]\\
&\quad+\Pr[P(\boldsymbol{\vec P})\notin\operatorname{Fake}(\boldsymbol\gamma,\mathbf E)\vee(\boldsymbol\gamma,\mathbf E)\notin\Xi]-3\delta\\
&\ge\Pr[P(\boldsymbol{\vec P})\notin\operatorname{Fake}(\boldsymbol\gamma,\mathbf E)\cap\operatorname{Fake}(\alpha(\boldsymbol\gamma,\mathbf E))\cap\operatorname{Fake}(\beta(\boldsymbol\gamma,\mathbf E))\vee(\boldsymbol\gamma,\mathbf E)\notin\Xi]-3\delta\\
&=1-3\delta .
\end{align*}
The last equality follows by the difference condition. Hence, $3(1-p)\ge1-3\delta$, thus $p\le2/3+\delta$.
\end{proof}
\begin{authorclaim}{Lemma 26}
For every 0-instance $\vec x$ of $\mathrm{UDISJ}_{k,2^{\ell-k}-1}$, for some $\delta<1/3-1/100$ there is a $(1-\delta)$-symmetry randomness space.
\end{authorclaim}
Lemma26 (proved in Section5.3) and Claim25 imply for a constant $\varepsilon>0$ that
\[
\Pr[P(\mathbf P_1,\ldots,\mathbf P_k)\notin\operatorname{Fake}(\boldsymbol\gamma,\mathbf E)]\ge\varepsilon+1/100,
\quad
\Pr[P(\mathbf P_1,\ldots,\mathbf P_k)\notin\operatorname{Fake}(\boldsymbol\gamma,\mathbf E)\cup\{\bot\}]\ge\varepsilon .
\]
Then, for $N=O(\log(1/\varepsilon))$, $T$ gives a correct answer for every 0-instance with probability at least $2/3$.
\end{proof}

\subsection*{5.3 Constructions of $\Xi,\alpha,\beta$}
\begin{proof}[Proof of Lemma 26]
Assume that $x_1,\ldots,x_k$ is a 0-instance $\mathrm{UDISJ}_{k,2^{\ell-k}-1}$. Let $i_0\in[2^{\ell-k}-1]$ be such that $x_{1,i_0}=x_{2,i_0}=\cdots=x_{k,i_0}=1$. Hereinafter $\gamma$ denotes a bijection from $[2^{\ell-k}-1]\cup\{*\}$ to $\{0,1\}^{\ell-k}$, $E$ denotes an invertible matrix from $\mathbb F_2^{\ell\times\ell}$ and $e_1,e_2,\ldots,e_\ell$ denote rows of $E$. Before presenting constructions of $\Xi,\alpha$, and $\beta$ we explain how we are going to establish symmetry and difference properties from Definition24. For every $s\in\{0,1\}^{\ell-k}$ and $b\in\{0,1\}$ we introduce the following notation:
\[
R(s,b,E):=\{(s,b,z)E\mid z\in\mathbb F_2^{k-1}\}.
\]
\begin{authorclaim}{Claim 27}
The last $2^{k+1}$ rows of $C_{\vec x}(\gamma,E)$ form $R(\gamma(*),0,E)\cup R(\gamma(*),1,E)$, each occurring twice. This is $\operatorname{Fake}(\gamma,E)$.
The rows with indices $[2^k(i-1)+1;2^ki]$ for $i\in[2^{\ell-k-1}]\setminus\{i_0\}$ are $R(\gamma(i),0,E)\cup R(\gamma(i),1,E)$, each occurring once.
The rows with indices $[2^k(i_0-1)+1;2^ki_0]$ are $R(\gamma(i_0),0,E)$, each occurring twice.
\end{authorclaim}
\begin{authorclaim}{Claim 28}
$R(s,b,E)=(\sum_{j=1}^{\ell-k}s_je_j+be_{\ell-k+1})+\operatorname{Span}(e_{\ell-k+2},\ldots,e_\ell)$.
\end{authorclaim}
\begin{proof}
\begin{align*}
R(s,b,E)&=\{(s,b,z)(e_1,\ldots,e_\ell)^T\mid z\in\mathbb F_2^{k-1}\}\\
&=\left\{\sum_{j=1}^{\ell-k}s_je_j+be_{\ell-k+1}+\sum_{i=1}^{k-1}z_ie_{\ell-k+1+i}\mid z\in\mathbb F_2^{k-1}\right\}\\
&=\left(\sum_{j=1}^{\ell-k}s_je_j+be_{\ell-k+1}\right)+\operatorname{Span}(e_{\ell-k+2},\ldots,e_\ell).
\end{align*}
\end{proof}
\begin{authorclaim}{Claim 29}
$|R(s,b,E)|=2^{k-1}$.
\end{authorclaim}
\begin{proof}
By Claim28, $|R(s,b,E)|=|(\sum_{j=1}^{\ell-k}s_je_j+be_{\ell-k+1})+\operatorname{Span}(e_{\ell-k+2},\ldots,e_\ell)|=|\operatorname{Span}(e_{\ell-k+2},\ldots,e_\ell)|=2^{k-1}$.
\end{proof}
\begin{authorclaim}{Claim 30}
The sets $R(s,b,E)$ for $s\in\{0,1\}^{\ell-k},b\in\{0,1\}$ are disjoint.
\end{authorclaim}
\begin{proof}
Consider two vectors $u\in R(s,b,E)$ and $v\in R(s',b',E)$ such that $(s,b)\ne(s',b')$. Then, by Claim28, $u$ and $v$ have different coordinates in the basis $e_1,e_2,\ldots,e_\ell$, hence $u\ne v$.
\end{proof}
\begin{authorclaim}{Claim 31}
For bijections $\gamma,\gamma'$ and invertible $E,E'$, if
\begin{align*}
&R(\gamma(i_0),0,E)\cup R(\gamma(*),0,E)\cup R(\gamma(*),1,E)\\
&\qquad=R(\gamma'(i_0),0,E')\cup R(\gamma'(*),0,E')\cup R(\gamma'(*),1,E'),\\
&R(\gamma(i_0),1,E)=R(\gamma'(i_0),1,E'),
\end{align*}
then $C_{\vec x}(\gamma,E)$ and $C_{\vec x}(\gamma',E')$ differ only by a row permutation.
\end{authorclaim}
\begin{proof}
By Claim27, rows from $R(\gamma(i_0),1,E)$ do not appear in $C_{\vec x}(\gamma,E)$, rows from $R(\gamma(i_0),0,E)\cup R(\gamma(*),0,E)\cup R(\gamma(*),1,E)$ appear in $C_{\vec x}(\gamma,E)$ exactly twice. The matrix $C_{\vec x}(\gamma,E)$ has $2^\ell+2^k$ rows. All rows of $C_{\vec x}(\gamma,E)$ that are not in $R(\gamma(i_0),1,E)\cup R(\gamma(*),0,E)\cup R(\gamma(*),1,E)$, by Claim27, appear in $C_{\vec x}(\gamma,E)$ exactly once.
By Claims29 and30, $|R(\gamma(i_0),0,E)\cup R(\gamma(*),0,E)\cup R(\gamma(*),1,E)|=3\cdot2^{k-1}$, hence, the number of rows of $C_{\vec x}(\gamma,E)$ that are not in $R(\gamma(i_0),1,E)\cup R(\gamma(*),0,E)\cup R(\gamma(*),1,E)$ equals $2^\ell-2^{k+1}$. By Claims29 and30, the number of $\ell$-bit strings not from $R(\gamma(i_0),1,E)\cup R(\gamma(i_0),0,E)\cup R(\gamma(*),0,E)\cup R(\gamma(*),1,E)$ is also $2^\ell-2^{k+1}$. Hence, all rows from $\{0,1\}^\ell\setminus(R(\gamma(i_0),0,E)\cup R(\gamma(*),0,E)\cup(\gamma(*),1,E)\cup R(\gamma(i_0),1,E))$ appear in $C_{\vec x}(\gamma,E)$ exactly once. Thus, matrices $C_{\vec x}(\gamma,E)$ and $C_{\vec x}(\gamma',E')$ have the same set of rows and each row appears the same number of times in each of these matrices.
\end{proof}
Write $\alpha(\gamma,E)=(\gamma_\alpha,E_\alpha)$ and $\beta(\gamma,E)=(\gamma_\beta,E_\beta)$. We construct them so that
\begin{equation}
\begin{aligned}
R(\gamma(i_0),1,E)&=R(\gamma_\alpha(i_0),1,E_\alpha)=R(\gamma_\beta(i_0),1,E_\beta),\\
R(\gamma(i_0),0,E)&=R(\gamma_\alpha(*),0,E_\alpha)=R(\gamma_\beta(*),0,E_\beta),\\
R(\gamma(*),1,E)&=R(\gamma_\alpha(*),1,E_\alpha)=R(\gamma_\beta(i_0),0,E_\beta),\\
R(\gamma(*),0,E)&=R(\gamma_\alpha(i_0),0,E_\alpha)=R(\gamma_\beta(*),1,E_\beta).
\end{aligned}\tag{2}
\end{equation}
Notice that by Claim31, equations(2) imply the symmetry property. Equations(2) also imply the difference property. Indeed, the three Fake sets are, respectively,
\begin{align*}
& R(\gamma(*),1,E)\cup R(\gamma(*),0,E),\\
& R(\gamma(*),1,E)\cup R(\gamma(i_0),0,E),\\
& R(\gamma(*),0,E)\cup R(\gamma(i_0),0,E),
\end{align*}
Hence, by Claim30, $\operatorname{Fake}(\gamma,E)\cap\operatorname{Fake}(\gamma_\alpha,E_\alpha)\cap\operatorname{Fake}(\gamma_\beta,E_\beta)=\varnothing$.
It remains to construct $\Xi,\alpha,\beta$ satisfying (2) and $\Pr[(\boldsymbol\gamma,\mathbf E)\in\Xi]>2/3+1/100$.
\paragraph{Definition of $\Xi$}
$(\gamma,E)\in\Xi$ iff there exist $m,n\in[\ell-k]$ with $\gamma(*)_m=1,\gamma(i_0)_m=0$ and $\gamma(*)_n=0,\gamma(i_0)_n=1$, i.e. $\gamma(*)$ and $\gamma(i_0)$ are coordinatewise incomparable.
$\boldsymbol\gamma(i_0)$ and $\boldsymbol\gamma(*)$ are uniform distinct elements. For independent uniform random subsets $\mathbf S,\mathbf T\subseteq[\ell-k]$,
\begin{align*}
\Pr[(\boldsymbol\gamma,\mathbf E)\in\Xi]
&=1-\Pr[\boldsymbol\gamma(i_0)\le\boldsymbol\gamma(*)\vee\boldsymbol\gamma(*)\le\boldsymbol\gamma(i_0)]\\
&\ge1-2\Pr[\boldsymbol\gamma(i_0)\le\boldsymbol\gamma(*)]\\
&=1-2\Pr[\mathbf S\subseteq\mathbf T\mid\mathbf S\ne\mathbf T]\\
&\ge1-2\Pr[\mathbf S\subseteq\mathbf T]\\
&=1-2\prod_{j=1}^{\ell-k}(1-\Pr[j\in\mathbf S\wedge j\notin\mathbf T])\\
&=1-2(3/4)^{\ell-k}>2/3+1/100\qquad(\ell-k\ge7).
\end{align*}
Hence largeness holds.
\paragraph{Construction of $\alpha$}
Let $E_\alpha$ have rows
\[
(e'_1,\ldots,e'_\ell)=
(e_1,\ldots,e_{\ell-k},e_{\ell-k+1}+\sum_{j=1}^{\ell-k}(\gamma(i_0)_j+\gamma(*)_j)e_j,e_{\ell-k+2},\ldots,e_\ell),
\]
and $\gamma_\alpha(i)=\gamma(*)$ for $i=i_0$, $\gamma(i_0)$ for $i=*$, and $\gamma(i)$ otherwise.
\begin{authorclaim}{Claim 32}
$\alpha:\Xi\to\Xi$ is bijective.
\end{authorclaim}
\begin{proof}
Notice that rows of $E'$ form a basis since $\sum_{j=1}^{\ell-k}(\gamma(i_0)_j+\gamma(*)_j)e_j\in\operatorname{Span}(e_1,\ldots,e_{\ell-k})$. The mapping $\gamma\mapsto\gamma_\alpha$ is bijective since it just swaps $\gamma(i_0)$ and $\gamma(*)$. Since the condition on $\gamma(i_0)$ and $\gamma(*)$ does not change after application of $\alpha$, we get that $\alpha(\Xi)\subseteq\Xi$. Notice that
\[
\sum_{j=1}^{\ell-k}(\gamma(i_0)_j+\gamma(*)_j)e_j
=\sum_{j=1}^{\ell-k}(\gamma_\alpha(i_0)_j+\gamma_\alpha(*)_j)e'_j,
\]
hence $\alpha(\gamma_\alpha,E_\alpha)=(\gamma,E)$, hence $\alpha$ is bijective.
\end{proof}
\begin{authorclaim}{Claim 33}
For all $(\gamma,E)\in\Xi$,
\begin{align*}
R(\gamma_\alpha(i_0),1,E_\alpha)&=R(\gamma(i_0),1,E),\\
R(\gamma_\alpha(i_0),0,E_\alpha)&=R(\gamma(*),0,E),\\
R(\gamma_\alpha(*),0,E_\alpha)&=R(\gamma(i_0),0,E),\\
R(\gamma_\alpha(*),1,E_\alpha)&=R(\gamma(*),1,E).
\end{align*}
\end{authorclaim}
\begin{proof}
Use Claim28 and $S:=\operatorname{Span}(e_{\ell-k+2},\ldots,e_\ell)=\operatorname{Span}(e'_{\ell-k+2},\ldots,e'_\ell)$:
\begin{align*}
R(\gamma_\alpha(i_0),1,E_\alpha)
&=(\sum_j\gamma_\alpha(i_0)_je'_j+e'_{\ell-k+1})+S\\
&=(\sum_j\gamma(*)_je'_j+e'_{\ell-k+1})+S
=(\sum_j\gamma(i_0)_je_j+e_{\ell-k+1})+S
=R(\gamma(i_0),1,E),\\
R(\gamma_\alpha(i_0),0,E_\alpha)
&=(\sum_j\gamma_\alpha(i_0)_je'_j)+S=(\sum_j\gamma(*)_je_j)+S=R(\gamma(*),0,E),\\
R(\gamma_\alpha(*),0,E_\alpha)
&=(\sum_j\gamma_\alpha(*)_je'_j)+S=(\sum_j\gamma(i_0)_je_j)+S=R(\gamma(i_0),0,E),\\
R(\gamma_\alpha(*),1,E_\alpha)
&=(\sum_j\gamma_\alpha(*)_je'_j+e'_{\ell-k+1})+S\\
&=(\sum_j\gamma(i_0)_je'_j+e'_{\ell-k+1})+S
=(\sum_j\gamma(*)_je_j+e_{\ell-k+1})+S=R(\gamma(*),1,E).
\end{align*}
Here each sum is over $j=1,\ldots,\ell-k$.
\end{proof}
\paragraph{Construction of $\beta$}
Set $\gamma_\beta=\gamma_\alpha$. Define $j_{\min}=\min\{j\in[\ell-k]:\gamma(*)_j=1\wedge\gamma(i_0)_j=0\}$, which exists on $\Xi$. Set $E_\beta=(e''_1,\ldots,e''_\ell)$ with
\[
e''_j=\begin{cases}
e_j,&j\notin\{j_{\min},\ell-k+1\},\\
\sum_{i=1}^{\ell-k}(\gamma(*)_i+\gamma(i_0)_i)e_i,&j=\ell-k+1,\\
e_{j_{\min}}+e_{\ell-k+1},&j=j_{\min}.
\end{cases}
\]
\begin{authorclaim}{Claim 34}
$\beta:\Xi\to\Xi$ is bijective.
\end{authorclaim}
\begin{proof}
Let us verify that $\beta$ is injective. Given $\gamma_\beta$ we can easily recover $\gamma$, hence we can recover $j_{\min}$ as well. Then
\begin{align*}
\sum_{i=1}^{\ell-k}(\gamma(i_0)_i+\gamma(*)_i)e''_i+e''_{\ell-k+1}
&=\sum_{i\in[\ell-k]\setminus\{j_{\min}\}}(\gamma(i_0)_i+\gamma(*)_i)e_i
+e_{j_{\min}}+e_{\ell-k+1}+e''_{\ell-k+1}\\
&=e_{\ell-k+1}+
\underbrace{\sum_{i\in[\ell-k]\setminus\{j_{\min}\}}(\gamma(i_0)_i+\gamma(*)_i)e_i+e_{j_{\min}}}_{e''_{\ell-k+1}}+e''_{\ell-k+1}\\
&=e_{\ell-k+1}.
\end{align*}
Thus, we can uniquely recover $e_{\ell-k+1}$ and, hence, also recover $e_{j_{\min}}=e''_{j_{\min}}+e_{\ell-k+1}$; for $j\in[\ell]\setminus\{j_{\min},\ell-k+1\}$, $e_j=e''_j$. Hence, $\beta$ is injective. Notice that since we represent $e_1,\ldots,e_\ell$ as linear combinations of $e''_1,\ldots,e''_\ell$, then $e''_1,\ldots,e''_\ell$ is a basis, hence the matrix $E_\beta$ is invertible. Thus, we verify that $\beta(\Xi)\subseteq\Xi$ and $\beta$ is injective, hence $\beta$ is bijective.
\end{proof}
\begin{authorclaim}{Claim 35}
For all $(\gamma,E)\in\Xi$,
\begin{align*}
R(\gamma_\beta(i_0),1,E_\beta)&=R(\gamma(i_0),1,E),\\
R(\gamma_\beta(i_0),0,E_\beta)&=R(\gamma(*),1,E),\\
R(\gamma_\beta(*),0,E_\beta)&=R(\gamma(i_0),0,E),\\
R(\gamma_\beta(*),1,E_\beta)&=R(\gamma(*),0,E).
\end{align*}
\end{authorclaim}
\begin{proof}
Put $S=\operatorname{Span}(e_{\ell-k+2},\ldots,e_\ell)=\operatorname{Span}(e''_{\ell-k+2},\ldots,e''_\ell)$. Recall $\gamma(*)_{j_{\min}}=1$, $\gamma(i_0)_{j_{\min}}=0$. Sums below run from1 to $\ell-k$:
\begin{align*}
R(\gamma_\beta(i_0),1,E_\beta)
&=\sum_i\gamma_\beta(i_0)_ie''_i+e''_{\ell-k+1}+S\\
&=\sum_i\gamma(*)_ie_i+e_{\ell-k+1}+e''_{\ell-k+1}+S\\
&=\sum_i\gamma(*)_ie_i+e_{\ell-k+1}+\sum_i(\gamma(*)_i+\gamma(i_0)_i)e_i+S\\
&=\sum_i\gamma(i_0)_ie_i+e_{\ell-k+1}+S=R(\gamma(i_0),1,E),\\
R(\gamma_\beta(i_0),0,E_\beta)
&=\sum_i\gamma_\beta(i_0)_ie''_i+S
=\sum_i\gamma(*)_ie_i+e_{\ell-k+1}+S=R(\gamma(*),1,E),\\
R(\gamma_\beta(*),0,E_\beta)
&=\sum_i\gamma_\beta(*)_ie''_i+S=\sum_i\gamma(i_0)_ie_i+S=R(\gamma(i_0),0,E),\\
R(\gamma_\beta(*),1,E_\beta)
&=\sum_i\gamma_\beta(*)_ie''_i+e''_{\ell-k+1}+S
=\sum_i\gamma(i_0)_ie_i+e''_{\ell-k+1}+S
=\sum_i\gamma(*)_ie_i+S=R(\gamma(*),0,E).
\end{align*}
\end{proof}
Claims33 and35 imply equations(2).
\end{proof}

\subsection*{5.4 Proof of Lemma 21}
To prove Lemma21 it is sufficient to prove:
\begin{authorclaim}{Proposition 36}
There exist matrices $T_1,\ldots,T_k\in\mathbb F_2^{2^k\times k}$ such that (1) for $\alpha_i\in\{0,1\}$, $\sum_i\alpha_iT_i=0$ iff all $\alpha_i=0$, i.e. they are linearly independent; (2) every nonzero $M\in\operatorname{Span}(T_1,\ldots,T_k)$ satisfies $M+K_k\in S_k$.
\end{authorclaim}
\begin{proof}[Proof of Lemma21]
Let for $i\in\{1,\ldots,k-1\}$, $A_i(0)=T_i$ and $A_i(1)$ be the zero matrix. Let $A_k(0)=K_k+T_k$, $A_k(1)=K_k$. For each $b_1,\ldots,b_k\in\{0,1\}$, $\sum_{i=1}^kA_i(b_i)=\sum_{i=1}^k(1-b_i)T_i+K_k$. Then $\sum_{i=1}^kA_i(1)=K_k$, and if for at least one $i\in[k]$, $b_i\ne1$, then by the first condition of Proposition36, $\sum_{i=1}^k(1-b_i)T_i$ differs from zero, thus by the second condition of Proposition36, $\sum_{i=1}^kA_i(b_i)\in S_k$.
\end{proof}
\begin{proof}[Proof of Proposition36]
Let us prove the proposition by induction on $k$. We are going to prove a stronger statement: namely, we additionally require that for arbitrary non-zero matrix $M\in\operatorname{Span}(T_1,\ldots,T_k)$ the set of even-indexed rows of $M+K_k\in S_k$ coincide with the set of odd-indexed rows of this matrix with all bits flipped.
The base case $k=1$. $T_1=\binom01$, and $K_1=\binom00$. It is easy to verify that all conditions hold.
Induction step from $k$ to $k+1$. Notice that
\[
K_{k+1}=\begin{pmatrix}K_k&0_{2^k\times1}\\K_k&1_{2^k\times1}\end{pmatrix}.
\]
Let $T_1,\ldots,T_k$ be the matrices from induction hypothesis for $k$. Then define
\[
T'_i=\begin{pmatrix}T_i&0_{2^k\times1}\\T_i&0_{2^k\times1}\end{pmatrix}
\]
for $i\in[k]$ and
\[
T'_{k+1}=\begin{pmatrix}0_{2^k\times k}&z_0\\1_{2^k\times k}&z_1\end{pmatrix},
\]
where $z_0=(0,1,0,1,\ldots,0,1)^T\in\{0,1\}^{2^k\times1}$, and $z_1=(1,0,1,0,\ldots,1,0)^T\in\{0,1\}^{2^k\times1}$.
Let us verify that all conditions hold. First we show that the matrices $T'_1,T'_2,\ldots,T'_{k+1}$ are linearly independent. Matrices $T'_1,T'_2,\ldots,T'_k$ are linearly independent since they contain linearly independent blocks $T_1,\ldots,T_k$. The matrix $T'_{k+1}$ does not belong to $\operatorname{Span}(T'_1,\ldots,T'_k)$, since the last column of $T'_{k+1}$ is non-zero, but the last columns of all $T'_1,\ldots,T'_k$ are zeros.
Let us check that for any non-zero matrix $M\in\operatorname{Span}(T'_1,\ldots,T'_k,T'_{k+1})$, the condition $M+K_{k+1}\in S_{k+1}$ holds and the set of even-indexed rows of $M+K_{k+1}$ coincide with the set of odd-indexed rows of this matrix with all bits flipped. Let us analyze the cases:
\begin{enumerate}
\item Let $M$ be a non-zero matrix from $\operatorname{Span}(T'_1,\ldots,T'_k)$. Then, $M$ has form
\[
\begin{pmatrix}M'&0_{2^k\times1}\\M'&0_{2^k\times1}\end{pmatrix},
\]
where $M'$ is a non-zero matrix from $\operatorname{Span}(T_1,\ldots,T_k)$, thus $M'+K_k\in S_k$. Then
\[
M+K_{k+1}=\begin{pmatrix}M'+K_k&0_{2^k\times1}\\M'+K_k&1_{2^k\times1}\end{pmatrix};
\]
it follows from the induction hypothesis that all rows of this matrix are distinct, i.e. $M+K_{k+1}\in S_{k+1}$. In order to verify that the set of even-indexed rows of this matrix coincide with the set of odd-indexed rows with all bits flipped, observe that by induction hypothesis the first $2^{k-1}$ even-indexed rows of $M+K_{k+1}$ coincide with the last $2^{k-1}$ odd-indexed rows of $M+K_{k+1}$ with all bits flipped, and the first $2^{k-1}$ odd-indexed rows of $M+K_{k+1}$ coincide with the last $2^{k-1}$ even-indexed rows of $M+K_{k+1}$ with flipped bits.
\item $M=T'_{k+1}$, then
\[
M+K_{k+1}=\begin{pmatrix}K_k&z_0\\1_{2^k\times k}+K_k&z_0\end{pmatrix}.
\]
Let us show that all rows of this matrix are distinct. The first $2^k$ rows start with0 and are obtained by appending zeroes and ones to the rows of $K_k$ in the alternating order. Since for every pair of coinciding rows of $K_k$ they are adjacent, the first $2^k$ rows are distinct. The last $2^k$ rows start from one, so they differ from the first $2^k$ rows. The proof that they are distinct is the same as for the first $2^k$ rows. Observe that the $(2i-1)$th row of the matrix $M+K_{k+1}$ coincide with the $(2^k+2i)$th row of $M+K_{k+1}$ with flipped bits, and the $(2i)$th row of $M+K_{k+1}$ coincide with the $(2^k+2i-1)$th row of $M+K_{k+1}$ with flipped bits for $i\in[2^k]$.
\item $M=R+T'_{k+1}$, where $R$ is a non-zero matrix from $\operatorname{Span}(T'_1,\ldots,T'_k)$. Let $R$ have the form
\[
\begin{pmatrix}R'&0_{2^k\times1}\\R'&0_{2^k\times1}\end{pmatrix},
\]
where $R'$ is a non-zero matrix from $\operatorname{Span}(T_1,\ldots,T_k)$. Then
\[
M+K_{k+1}=R+T'_{k+1}+K_{k+1}
=\begin{pmatrix}R'+K_k&z_0\\1_{2^k\times k}+R'+K_k&z_0\end{pmatrix}.
\]
By the induction hypothesis, $R'+K_k\in S_k$ and its even-indexed rows coincide with its odd-indexed rows with flipped bits. Then, all even-indexed rows of $M+K_{k+1}$ end with0, the first $2^{k-1}$ of them are even-indexed rows of $R'+K_k$ with appended zero, and the last $2^{k-1}$ of them are even-indexed rows of $R'+K_k$ with all bits flipped and appended0. Then, by the induction hypothesis, the set of the former rows does not intersect with the set of the latter rows, therefore they are all distinct. By the same argument, all the rows of $M+K_{k+1}$ that end with1 are distinct. Thus, $M+K_{k+1}\in S_{k+1}$.
Let us verify that the set of even-indexed rows of this matrix coincide with the set of odd-indexed rows of this matrix with all bits flipped. Observe that if the $i$th row of $R'+K_k$ coincides with the $j$th row of $R'+K_k$ with flipped bits, then the $i$th row of $M+K_{k+1}$ coincides with its $j$th row with flipped bits, and the $(2^k+i)$th row of $M+K_{k+1}$ coincides with its $(2^k+j)$th row with all bits flipped. The required property follows from the induction hypothesis.
\end{enumerate}
\end{proof}
\subsection*{5.5 Corollaries}
\begin{authorclaim}{Corollary37}
If $k+7\le\ell$, the size of any semantic $\mathrm{Res}(\mathrm{PC}_{k-1})$ tree-like refutation of $\mathrm{BPHP}_{2^\ell}^{2^\ell+2^k}$ is at least $2^{\Omega(2^{\ell/2}/(k2^{3k/2}))}$. For $k=2$, every tree-like semantic $\mathrm{Res}(\oplus)$ refutation of $\mathrm{BPHP}_{2^\ell}^{2^\ell+4}$ has size at least $2^{\Omega(2^\ell)}$.
\end{authorclaim}
\begin{proof}
Follows from Theorem3 and Lemma1.
\end{proof}
\begin{authorclaim}{Corollary38}
Let $2\le k\le\ell-7$ and $S$ be the minimal size of tree-like refutation of $\phi=\mathrm{BPHP}_{2^\ell}^{2^\ell+2^k}\circ\oplus_k$ in the semantic proof system $\mathrm T^{\mathrm{cc}}(k,c)$. Then
$\log S\log\log S\ge c\cdot\Omega(2^{\ell/2}/(k2^{3k/2}))$.
For $k=2$, $\log S\log\log S\ge c\cdot\Omega(2^\ell)$.
\end{authorclaim}
\begin{proof}
By Lemma9, $R^{1/3}_{\mathrm{pub}}(\operatorname{Search}(\phi))=O(\log S\log\log S/c)$. Also
\[
R^{1/3}_{\mathrm{pub}}(\operatorname{Search}(\mathrm{BPHP}_{2^\ell}^{2^\ell+2^k}\circ\oplus_k))
\ge R^{1/3}_{\mathrm{pub}}(\oplus_k\operatorname{Search}(\mathrm{BPHP}_{2^\ell}^{2^\ell+2^k})).
\]
Now apply Theorem3.
\end{proof}

\section*{6 Bit pigeonhole principle}
\subsection*{6.1 Reduction from $\mathrm{BPHP}\circ\oplus_k$ to $\mathrm{BPHP}$}
Let $T\subseteq X_1\times\cdots\times X_k\times Y$ and $S\subseteq Z_1\times\cdots\times Z_k\times W$ be relations. We say $S$ is many-one reducible to $T$ if there are maps $f_i:X_i\to Z_i$ and $g:W\to Y$ such that if $(f_1(x_1),\ldots,f_k(x_k),y)\in T$ then $(x_1,\ldots,x_k,g(y))\in S$.
\begin{authorclaim}{Lemma39}
If $S$ is many-one reducible to $T$, then $R^{\mathrm{pub}}_{1/3}(S)\le R^{\mathrm{pub}}_{1/3}(T)$.
\end{authorclaim}
\begin{proof}
The $i$th party computes $f(x_j)$ for all $j\in[k]\setminus\{i\}$ and then all parties run the optimal protocol for $T$. As soon as all the parties learn an answer $y$ they compute $g(y)$ without communication.
\end{proof}
Recall $\mathrm{BPHP}_{2^n}^M$ encodes $M$ different strings $s_1,\ldots,s_M\in\{0,1\}^n$. For $n=\ell k+r$, $0\le r<k$, partition each row $s_i=s_i^{(1)}\cdots s_i^{(k)}$, where $|s_i^{(t)}|=\ell+1$ for $t\le r$ and $\ell$ for $t>r$. Partition $\Pi_k$ puts all $s_1^{(t)},\ldots,s_M^{(t)}$ in part $t$.
$\operatorname{SearchPair}_{2^n}^M$ asks for distinct $i,j$ with $s_i=s_j$ from variables partitioned according to $\Pi_k$.
\begin{authorclaim}{Proposition40}
$\operatorname{SearchPair}_{2^n}^M$ is many-one reducible to $\operatorname{Search}(\mathrm{BPHP}_{2^n}^M)$ with partition $\Pi_k$.
\end{authorclaim}
\begin{proof}
The proof is straightforward.
\end{proof}
\begin{authorclaim}{Theorem41}
$\oplus_k\mathrm{BPHP}_{2^\ell}^m$ is many-one reducible to $\operatorname{SearchPair}_{2^{k\ell}}^{m2^{(k-1)\ell}}$.
\end{authorclaim}
\begin{proof}
Let us denote $M=m\cdot2^{(k-1)\ell}$. Consider a set
$Z=\{(y_1,y_2,\ldots,y_k)\in(\mathbb F_2^\ell)^k\mid\sum_iy_i=0\}$.
It is easy to see that $|Z|=2^{(k-1)\ell}$. Let $\phi$ be a bijection between $[M]$ and $Z\times[m]$.
Let for $i\in[m]$ and $t\in[k]$, $x_i^{(t)}$ denote the $i$th string of the $t$th party in the communication problem $\oplus_k\mathrm{BPHP}_{2^\ell}^m$. Let $x_i:=(x_i^{(1)},\ldots,x_i^{(k)})$.
For every $t\in[k]$ we define $f_t$ as follows: $f_t(x_1^{(t)},\ldots,x_m^{(t)})$ is a sequence of rows $r_1^{(t)},r_2^{(t)},\ldots,r_M^{(t)}$ such that for all $i\in[M]$, $r_i^{(t)}=z_t+x_j^{(t)}$, where $(z,j)=\phi(i)$ for all $z\in Z$ and $j\in[m]$ (recall that $z\in Z$ is divided on $k$ parts of equal lengths and $z_t$ denotes the $t$th part).
Let us construct the function $g$ from the definition of the reduction.
Let $q,w\in[M]$ and $q\ne w$. Assume that $\phi(q)=(z,j_1)$ and $\phi(w)=(z,j_2)$. We define $g(q,w):=(j_1,j_2)$.
Let us verify that $f_1,f_2,\ldots,f_k$ and $g$ define a reduction. Let $q,w\in M$ be a pair of different numbers such that the assignment $\alpha:=\{s_i\leftarrow r_i^{(1)}r_i^{(2)}\cdots r_i^{(k)}\mid i\in[M]\}$ satisfies $s_q=s_w$.
Assume that $g(q,w)=(j_1,j_2)$. We need to verify that $j_1\ne j_2$ and $\sum_{t=1}^kx_{j_1}^{(t)}=\sum_{t=1}^kx_{j_2}^{(t)}$.
Notice that under the assignment $\alpha$ the value of $s_q$ is $x_{j_1}+z$ and the value of $s_w$ is $x_{j_2}+y$, where $j_1,j_2\in[m]$ and $z,y\in Z$ such that $(z,j_1)=\phi(q)$ and $(y,j_2)=\phi(w)$. If $j_1=j_2$, then $x_{j_1}+z=x_{j_2}+y$ implies $z=y$. Since $\phi$ is a bijection, we get $q=w$. Thus, $j_1\ne j_2$.
For each $t\in[k]$, the following equality holds.
\begin{equation}
z_t+x_{j_1}^{(t)}=y_t+x_{j_2}^{(t)}.\tag{3}
\end{equation}
If we sum up equations(3) for all $t\in[k]$ and use that $y,z\in Z$, we get $\sum_{t=1}^kx_{j_1}^{(t)}=\sum_{t=1}^kx_{j_2}^{(t)}$. Hence, $(j_1,j_2)$ is a correct answer for $\oplus_k\mathrm{BPHP}_{2^\ell}^m$.
\end{proof}
The following proposition deals with the case, where the number of bits is not divisible by $k$.
\begin{authorclaim}{Proposition42}
For $n=k\ell+r$, $0\le r<k$, $M>2^{k\ell}$, $\operatorname{SearchPair}_{2^{k\ell}}^M$ is many-one reducible to $\operatorname{SearchPair}_{2^n}^{M2^r}$.
\end{authorclaim}
\begin{proof}
Let $x_1,x_2,\ldots,x_M$ be the input of $\operatorname{SearchPair}_{2^{k\ell}}^M$, let $x_j^{(t)}$ be the $t$th part of the row $x_j$ according to the partition $\Pi_k$. Given this input we construct an input for $\operatorname{SearchPair}_{2^n}^{M2^r}$.
Let $\tau$ be a bijection between $[M]\times\{0,1\}^r$ and $[M2^r]$.
For each $i\in[M]$ we construct $2^r$ rows $y_{\tau(i,\alpha)}$ one for each $\alpha\in\{0,1\}^r$. Let $\Pi_k$ partition a row $y_{\tau(i,\alpha)}$ into the following parts: $y_{\tau(i,\alpha)}^{(1)}y_{\tau(i,\alpha)}^{(2)}\cdots y_{\tau(i,\alpha)}^{(k)}$. Let
\[
y_{\tau(i,\alpha)}^{(t)}
=\begin{cases}x_i^{(t)},&t>r,\\x_i^{(t)}\alpha_t,&0\le t\le r.\end{cases}
\]
Now we can define the function $f_t(x_1^{(t)},\ldots,x_M^{(t)})$ as $y_{\tau(i,\alpha)}^{(t)}$ for each $i\in M$ and $\alpha\in\{0,1\}^r$ and $t\in[k]$. Observe that for each $i\in[M]$ the rows $y_{i,\alpha}$ for $\alpha\in\{0,1\}^r$ are distinct. That allows us to define the function $g$ as $g(\tau(i_1,\alpha_1),\tau(i_2,\alpha_2))=(i_1,i_2)$. All the required properties can be easily verified.
\end{proof}
\begin{authorclaim}{Theorem4}
Let $M=2^n+2^{k+n-\lfloor n/k\rfloor}$ and $n\ge k(k+7)$. With partition $\Pi_k$,
\[
R^{\mathrm{pub}}_{1/3}(\operatorname{Search}(\mathrm{BPHP}_{2^n}^M))
=\Omega(2^{n/(2k)-3k/2}/k).
\]
For $k=2$ the stronger bound is $\Omega(2^{n/2})$.
\end{authorclaim}
\begin{proof}
Let $\ell=\lfloor n/k\rfloor$, $r=n-\ell k$. By Proposition40, Proposition42, Theorem41 and Corollary19, respectively,
\begin{align*}
R^{\mathrm{pub}}_{1/3}(\operatorname{Search}(\mathrm{BPHP}_{2^n}^M))
&=R^{\mathrm{pub}}_{1/3}(\operatorname{Search}(\mathrm{BPHP}_{2^n}^{(2^k+2^\ell)2^{(k-1)\ell+r}}))\\
&\ge R^{\mathrm{pub}}_{1/3}(\operatorname{SearchPair}_{2^n}^{(2^k+2^\ell)2^{(k-1)\ell+r}})\\
&\ge R^{\mathrm{pub}}_{1/3}(\operatorname{SearchPair}_{2^{k\ell}}^{(2^k+2^\ell)2^{(k-1)\ell}})\\
&\ge R^{\mathrm{pub}}_{1/3}(\oplus_k\mathrm{BPHP}_{2^\ell}^{2^\ell+2^k})\\
&=\Omega(2^{\ell/2-3k/2}/k)=\Omega(2^{n/(2k)-3k/2}/k).
\end{align*}
The case of $k=2$ can be treated in the same way, the only difference is in the application of Corollary19.
\end{proof}
\subsection*{6.2 Upper bound for communication complexity}
\begin{authorclaim}{Proposition5}
For $M>2^n$ and $k\in\{2,\ldots,n\}$ there is a deterministic NOF protocol for $\operatorname{Search}(\mathrm{BPHP}_{2^n}^M)$ with partition $\Pi_k$ transmitting $O(2^{\lceil n/k\rceil}\log M)$ bits.
\end{authorclaim}
\begin{proof}
The protocol is going to have only two active parties: the second party, which we call Alice, and the first party, which we call Bob. We are going to use that Alice can see the variables $s_1^{(1)},\ldots,s_M^{(1)}$ and that Bob can see all other variables.
Let us denote $\bar s_i=s_i^{(1)}s_i^{(2)}s_i^{(3)}\cdots s_i^{(k)}\in\{0,1\}^{n-\lceil n/k\rceil}$ the bits Bob sees in the $i$th line for $i\in[M]$. Bob finds a value $\alpha\in\{0,1\}^{n-\lceil n/k\rceil}$ such that the size of the set $S_\alpha=\{i\in[M]\mid\bar s_i^{(1)}=\alpha\}$ is larger than $2^{\lceil n/k\rceil}$. Such $\alpha$ exists since $M>2^n$. Bob then picks an arbitrary subset $S'$ of $S_\alpha$ of size $2^{\lceil n/k\rceil}+1$ and sends the description of $S'$ to Alice using $(2^{\lceil n/k\rceil}+1)\cdot\lceil\log_2M\rceil$ bits. Then, by the pigeonhole principle there exists $i\ne j\in S'$ such that $s_i^{(1)}=s_j^{(1)}$. Alice and Bob then spend $O(\log M+n)$ bits transmitting indices $i$ and $j$ and all the values of the $i$th and $j$th lines to each other. Both of them then find the falsified clause of $\mathrm{BPHP}_{2^n}^M$ with no communication because it only depends on variables $s_i$ and $s_j$ and broadcast its description to all of the parties using an additional $O(n+\log M)$ bits.
\end{proof}
For $k=2$ this coincides with the lower bound up to a logarithmic factor. For larger $k$ they are polynomially related; the dependence on $k$ is a genuine phenomenon.
\subsection*{6.3 Short $\mathrm{Th}(\log n)$ proof of $\mathrm{BPHP}_n^m$}
This gives a short tree-like $\mathrm{Th}(\log n)$ refutation, similar to the conversion of unary PHP resolution proofs to $\mathrm{Res}(\log n)$ proofs of binary PHP in [5].
\begin{authorclaim}{Proposition43}
If there exists a tree-like $\mathrm{Th}(1)$ refutation of $\mathrm{PHP}_{2^\ell}^m$ of size $S$, there exists a tree-like $\mathrm{Th}(\ell)$ refutation of $\mathrm{BPHP}_{2^\ell}^m$ of size $O(S)$.
\end{authorclaim}
\begin{proof}
Let $p_{i,j}$ for $i\in[m]$ and $j\in[2^\ell]$ be a variable of $\mathrm{PHP}_{2^\ell}^m$ indicating that the $i$th pigeon flies to the $j$th hole. Let $s_{i,k}$ for $i\in[m]$, $k\in[\ell]$ be a variable of $\mathrm{BPHP}_{2^\ell}^m$ indicating the $\ell$th bit of the $i$th string $s_i$.
Let $Q_j(x_1,x_2,\ldots,x_\ell)$ for $j\in[2^k]$ be a multilinear polynomial over reals such that for all $a_1,a_2,\ldots,a_\ell\in\{0,1\}^\ell$, $Q_j(a_1,a_2,\ldots,a_\ell)=1$ if $(a_1,a_2,\ldots,a_\ell)=\operatorname{bin}_\ell(j-1)$ and $Q_j(a_1,a_2,\ldots,a_\ell)=0$ otherwise. We ma define $Q_j$ as follows
\[
Q_j(x_1,\ldots,x_\ell)=\prod_{k=1}^{\ell}(1-x_k+\alpha_\ell)
\]
for $i\in[m]$, $j\in[2^k]$, where $\alpha=\operatorname{bin}_\ell(j-1)$. By the construction $\deg(Q_j)=\ell$.
Let $P_{i,j}=Q_j(s_{i,1},s_{i,2},\ldots,s_{i,\ell})$.
Consider a tree-like $\mathrm{Th}(1)$-refutation of $\mathrm{PHP}_{2^\ell}^m$ of size $S$: $f_1\ge0,f_2\ge0,\ldots,f_S\ge0$, where $f_i$ are linear real polynomials over variables $p_{i,j}$ and $f_S\ge0$ is unsatisfiable on Boolean cube. For each of the inequalities on the following conditions hold:
\begin{enumerate}
\item[(a)] $f_i\ge0$ is semantically implied by $f_j\ge0$ and $f_k\ge0$ on the Boolean cube for $j,k<i$.
\item[(b)] $f_i$ is a linear representation of an axiom of $\mathrm{PHP}_{2^\ell}^m$;
\end{enumerate}
Let $F_i$ be a polynomial obtained of substitution $p_{j,k}:=P_{j,k}$ to $f_i$ for all $j\in[m]$; $k\in[2^\ell]$.
Consider a sequence of inequalities $F_1\ge0,\ldots,F_S\ge0$. Observe that $F_S\ge0$ is unsatisfiable on the Boolean cube since $P_{i,j}\in\{0,1\}$ on the Boolean cube. Let us verify that the sequence $F_1\ge0,\ldots,F_S\ge0$ may be extended to a correct tree-like $\mathrm{Th}(\ell)$ refutation of $\mathrm{BPHP}_{2^\ell}^m$:
\begin{enumerate}
\item[(a)] If $f_i\ge0$ is semantically implied by $f_j\ge0$ and $f_k\ge0$, then $F_i\ge0$ is also implied by $F_j\ge0$ and $F_k\ge0$, since $P_{i,j}$ is Boolean on the Boolean cube.
\item[(b)] If $f_i$ is a linear representation of a
\begin{itemize}
\item hole axiom then $f_i\ge0$ is equivalent to the function $(1-p_{a,b})+(1-p_{c,b})\ge1$ on $\{0,1\}^{\operatorname{Vars}(\mathrm{PHP}_{2^\ell}^m)}$ for $a,c\in[m]$, $b\in[2^\ell]$. Thus $F_i\ge0$ is also equivalent to $(1-P_{a,b})+(1-P_{c,b})\ge1$ on the Boolean cube. Observe that the restriction of $(1-P_{a,b})+(1-P_{c,b})\ge1$ to the Boolean cube coincides with the predicate $s_a\ne\operatorname{bin}_\ell(b)\vee s_c\ne\operatorname{bin}_\ell(b)$ which is an axiom of $\mathrm{BPHP}_{2^\ell}^m$.
\item pigeon axiom then $f_i\ge0$ is equivalent to $\sum_{j=1}^{2^\ell}p_{a,j}\ge1$ on the Boolean cube for some $a\in[m]$. Thus $F_i\ge0$ is equivalent to $\sum_{j=1}^{2^\ell}P_{a,j}\ge1$ on $\{0,1\}^{\operatorname{Vars}(\mathrm{BPHP}_{2^\ell}^m)}$. Observe that the latter inequality is identically true, since $P_{a,j}$ is equivalent to $s_a=\operatorname{bin}_\ell(j-1)$, so for exactly one value of $j\in[2^\ell]$, $P_{a,j}=1$. Since $F_i\ge0$ is identically true it can be semantically derived from two arbitrary axioms of $\mathrm{BPHP}_{2^\ell}^m$.
\end{itemize}
\end{enumerate}
It is easy to see that the size of the resulting refutation is at most $3S$.
\end{proof}
\begin{authorclaim}{Proposition44 ([4])}
For $m>n$ there is a tree-like Cutting Planes (a subsystem of $\mathrm{Th}(1)$) refutation of $\mathrm{PHP}_n^m$ of size $O(m^2n)$.
\end{authorclaim}
\begin{authorclaim}{Proposition6}
For $m>2^\ell$ there is a tree-like $\mathrm{Th}(\ell)$ refutation of $\mathrm{BPHP}_{2^\ell}^m$ of size $O(m^2 2^\ell)$.
\end{authorclaim}
\begin{proof}
Follows from Propositions43 and44.
\end{proof}

\section*{A Proof of Lemma 12}
\begin{authorclaim}{Lemma 12}
Let $T$ be a binary tree with $m$ vertices such that the $i$th vertex is labeled with $a_i\in\{0,1\}$ with the hereditary property: for each inner vertex $i$ with direct descendants $c_1$ and $c_2$, if $a_i=1$, then $a_{c_1}=1$ or $a_{c_2}=1$. We also assume that if $r$ is the root of $T$, then $a_r=1$. Assume that we have a one-sided bounded error oracle access to $a_i$ i.e. if we request a value of $a_i$ and $a_i=0$ we get 1 with probability at most $1/2$ and 0 with probability at least $1/2$; if $a_i=1$ we get 1 with probability 1. Then there exists an algorithm $A$ that with probability at least $2/3$ returns a leaf $\ell$ of $T$ with $a_\ell=1$ and makes $O(\log m)$ oracle queries to $a_1,\ldots,a_m$.
\end{authorclaim}
\begin{proof}[Proof of Lemma 12]
For a tree $F$ we denote by $|F|$ the number of nodes in $F$ and for a node $v$ of $F$ we denote by $\operatorname{Subtree}(F,v)$ the subtree of $F$ with root $v$. Let Oracle$(i)$ be the oracle function returning the correct value of $a_i$ with probability at least $9/10$. We can implement such a function using the majority vote of a constant number of initial oracle queries. Let $C$ be a constant; an appropriate value of $C$ we choose later. Consider Algorithm 2 on the following page.
We claim that at any iteration $T_i$ has the hereditary property. This is the case in the beginning and if $i$ decreases at some iteration, then the next $T_i$ was considered at an earlier iteration. Otherwise, the next $T_i$ is either a subtree of the current $T_i$ (in that case the hereditary property is clearly maintained), or is obtained by removal a subtree with 0-labeled root (here we use that the oracle has a one-sided error) from the previous $T_i$ (the hereditary property is also maintained in that case).
We first consider a variant of the algorithm that works infinitely long (i.e., $C=+\infty$) and compute the expected number of the first iteration such that $T_i$ consists of a single 1-labeled leaf of $T$. Notice that after the first such iteration the value of $T_i$ stays the same for all further iterations. We show that that the expected value is at most $C\log m$ for some constant $C$. Then by running the algorithm for $3C\lceil\log m\rceil$ iterations we obtain the required error probability by Markov's inequality.
\begin{authorclaim}{Algorithm 2: Search for 1-leaf}
\begin{algorithmic}
\State $T_0:=T$ \Comment{Initialize the tree}
\State $i:=0$
\For{$j:=1$ to $3C\lceil\log_{3/2}m\rceil$}
 \State $r:=$ root of $T_i$
 \If{Oracle$(r)=0$}
  \State $i:=\max\{0,i-1\}$ \Comment{Backtrack since the current tree may not contain a 1-leaf}
 \ElsIf{$|T_i|\ne1$}
  \State $v:=$ a centroid node of $T_i$
  \Statex \hspace{\algorithmicindent} $\triangleright$ i.e. such that $|\operatorname{Subtree}(T_i,v)|\in[\frac13|T_i|,\frac23|T_i|]$
  \If{Oracle$(v)=1$}
   \State $T_{i+1}:=\operatorname{Subtree}(T_i,v)$
  \Else
   \State $T_{i+1}:=T_i-\operatorname{Subtree}(T_i,v)$
   \Statex \hspace{\algorithmicindent} $\triangleright$ $T_{i+1}$ is obtained from $T_i$ by the deletion of $\operatorname{Subtree}(T_i,v)$
  \EndIf
  \State $i:=i+1$
 \EndIf
\EndFor
\State \Return the only node of $T_i$, if $|T_i|=1$
\end{algorithmic}
\end{authorclaim}
Let $T(j)$ denote the value of $T_i$ before the start of $j$th iteration, $i(j)$ denote $i$ at the start of $j$th iteration and $r(j)$ denote the root of $T(j)$. Notice that if $a_{r(j)}=1$, then for every $j'>j$, $T(j')$ is a subtree of $T(j)$, since the algorithm never backtracks if the true value of the roots label is 1. Hence, if $a_{r(j)}=a_{r(j')}=1$ for some $j<j'$, then $i(j)\le i(j')$.
Let us consider a sequence $j_1,j_2,j_3,\ldots$, where $j_1=0$, $j_s=\min\{j\mid a_{r(j)}=1\wedge j>j_{s-1}\wedge i(j)>i(j_{s-1})\}$, if such minimum exists.
Let us consider the iterations from $j_s$ till $j_{s+1}-1$. We consider the random variables $Y_{j_s},Y_{j_s+1},\ldots,Y_{j_{s+1}-1}$ corresponding to these iterations with the following properties:
\begin{itemize}
\item If $T(j)$ coincides with $T(j_s)$, then its root is labeled with 1. Then $Y_j=-1$ if the second oracle query returns the correct answer and $Y_j=1$ if the answer it incorrect. Notice that $\Pr[Y_j=-1]\ge9/10$.
\item If the root of $T(j)$ is labeled with zero, then $Y_j=-1$, if the first oracle query returns the correct answer (i.e. the algorithm backtracks). Otherwise, if $T(j)$ consists of a single node $Y_j=0$. Otherwise, if the root of $T(j+1)$ is labeled with 0, then $Y_j=1$. If it is labeled with 1, then $Y_j=-\infty$. Notice that $\Pr[Y_j\le-1]\ge9/10$.
\end{itemize}
Notice that, $j_{s+1}=j_s+\min\{k\mid\sum_{j=j_s}^{j_s+k-1}Y_j\le-1\}$. In order to estimate the expected value of $j_{s+1}-j_s$ we consider an auxiliary random variables $X_{j_s},X_{j_s+1},\ldots,X_{j_{s+1}-1}$, defined as
\[
X_j=\begin{cases}1,&\text{if }Y_j\ge0\\-1,&\text{if }Y_j<0.\end{cases}
\]
Notice then $\sum_{j=j_s}^{j_s+k-1}Y_j\le\sum_{j=j_s}^{j_s+k-1}X_j$. We can apply the following fact about random walks in a straight line to the random variables $X_j$:
\begin{authorclaim}{Theorem 45 (Section XII.2 of [18])}
Let $X_1,X_2,\ldots$ be a sequence of independent random variables that take value in $\{-1,1\}$. Assume that for all $i$, $\Pr[X_i=1]\le1/10$ and $\Pr[X_i=-1]\ge9/10$. Let $M$ be a random variable that equals the minimal natural number $k$ such that $\sum_{i=1}^kX_i=-1$. Then the expected value of $M$ is at most $C$, where $C\in\mathbb R$ is an absolute constant.
\end{authorclaim}
Fact 45 implies that $\mathbb E[j_{s+1}-j_s]\le C$. Then $\mathbb E[j_s]=\mathbb E[j_s-j_{s-1}+(j_{s-1}-j_{s-2})+\cdots+(j_2-j_1)+(j_1-j_0)]\le sC$. Thus, by Markov's inequality $\Pr[j_s\le3sC]\ge2/3$. Since $|T_{j_s}|\le(2/3)^s|T_{j_0}|$, the algorithm that runs for $3C\lceil\log_{3/2}m\rceil$ iterations terminates in a 1-labeled leaf with probability at least $2/3$.
\end{proof}

\begin{thebibliography}{33}
\bibitem{r1} Paul Beame, Toniann Pitassi, and Nathan Segerlind. Lower bounds for Lovász-Schrijver systems and beyond follow from multiparty communication complexity. SIAM J. Comput., 37(3):845–869, 2007.
\bibitem{r2} Paul Beame and Søren Riis. More on the relative strength of counting principles. In Paul Beam and Samuel R. Buss, editors, Proof Complexity and Feasible Arithmetics, Proceedings of a DIMACS Workshop, New Brunswick, New Jersey, USA, April 21-24, 1996, volume 39 of DIMACS Series in Discrete Mathematics and Theoretical Computer Science, pages 13–35. DIMACS/AMS, 1996. doi:10.1090/dimacs/039/02.
\bibitem{r3} Stephen A. Cook and Robert A. Reckhow. The relative efficiency of propositional proof systems. The Journal of Symbolic Logic, 44(1):36–50, 1979.
\bibitem{r4} William Cook, Collette R. Coullard, and Gy. Turán. On the complexity of cutting-plane proofs. Discrete Applied Mathematics, 18(1):25–38, 1987.
\bibitem{r5} Stefan S. Dantchev, Nicola Galesi, and Barnaby Martin. Resolution and the binary encoding of combinatorial principles. In Amir Shpilka, editor, 34th Computational Complexity Conference, CCC 2019, July 18-20, 2019, New Brunswick, NJ, USA, volume 137 of LIPIcs, pages 6:1–6:25. Schloss Dagstuhl - Leibniz-Zentrum für Informatik, 2019. doi:10.4230/LIPIcs.CCC.2019.6.
\bibitem{r6} Noah Fleming, Denis Pankratov, Toniann Pitassi, and Robert Robere. Random $\Theta(\log n)$CNFs are hard for cutting planes. In Chris Umans, editor, 58th IEEE Annual Symposium on Foundations of Computer Science, FOCS 2017, Berkeley, CA, USA, October 15-17, 2017, pages 109–120. IEEE Computer Society, 2017. doi:10.1109/FOCS.2017.19.
\bibitem{r7} Ankit Garg, Mika Göös, Pritish Kamath, and Dmitry Sokolov. Monotone circuit lower bounds from resolution. In Ilias Diakonikolas, David Kempe, and Monika Henzinger, editors, Proceedings of the 50th Annual ACM SIGACT Symposium on Theory of Computing, STOC 2018, Los Angeles, CA, USA, June 25-29, 2018, pages 902–911. ACM, 2018. doi:10.1145/ 3188745.3188838.
\bibitem{r8} Michal Garlík and Leszek Aleksander Kolodziejczyk. Some subsystems of constant-depth frege with parity. ACM Trans. Comput. Log., 19(4):29:1–29:34, 2018. doi:10.1145/3243126.
\bibitem{r9} Mika Göös and Toniann Pitassi. Communication lower bounds via critical block sensitivity. In Proceedings of the Forty-Sixth Annual ACM Symposium on Theory of Computing, STOC ’14, page 847–856, New York, NY, USA, 2014. Association for Computing Machinery. doi: 10.1145/2591796.2591838.
\bibitem{r10} Pavel Hrubes and Pavel Pudlák. Random formulas, monotone circuits, and interpolation. In Chris Umans, editor, 58th IEEE Annual Symposium on Foundations of Computer Science, FOCS 2017, Berkeley, CA, USA, October 15-17, 2017, pages 121–131. IEEE Computer Society, 2017. doi:10.1109/FOCS.2017.20.
\bibitem{r11} Pavel Hrubes and Pavel Pudlák. A note on monotone real circuits. Inf. Process. Lett., 131:15–19, 2018. doi:10.1016/j.ipl.2017.11.002.
\bibitem{r12} Trinh Huynh and Jakob Nordström. On the virtue of succinct proofs: amplifying communication complexity hardness to time-space trade-offs in proof complexity. In Howard J. Karloff and Toniann Pitassi, editors, Proceedings of the 44th Symposium on Theory of Computing Conference, STOC 2012, New York, NY, USA, May 19 - 22, 2012, pages 233–248. ACM, 2012. doi:10.1145/2213977.2214000.
\bibitem{r13} Russell Impagliazzo, Toniann Pitassi, and Alasdair Urquhart. Upper and lower bounds for tree-like cutting planes proofs. In Proceedings of the Ninth Annual Symposium on Logic in Computer Science (LICS ’94), Paris, France, July 4-7, 1994, pages 220–228. IEEE Computer Society, 1994. doi:10.1109/LICS.1994.316069.
\bibitem{r14} Dmitry Itsykson and Dmitry Sokolov. Lower bounds for splittings by linear combinations. In Erzsébet Csuhaj-Varjú, Martin Dietzfelbinger, and Zoltán Ésik, editors, Mathematical Foundations of Computer Science 2014 - 39th International Symposium, MFCS 2014, Budapest, Hungary, August 25-29, 2014. Proceedings, Part II, volume 8635 of Lecture Notes in Computer Science, pages 372–383. Springer, 2014. doi:10.1007/978-3-662-44465-8\_32.
\bibitem{r15} Dmitry Itsykson and Dmitry Sokolov. Resolution over linear equations modulo two. Annals of Pure and Applied Logic, 171(1), January 2020. doi:10.1016/j.apal.2019.102722.
\bibitem{r16} Bala Kalyanasundaram and Georg Schnitger. The probabilistic communication complexity of set intersection. SIAM J. Discret. Math., 5(4):545–557, 1992. doi:10.1137/0405044.
\bibitem{r17} Mauricio Karchmer and Avi Wigderson. Monotone circuits for connectivity require superlogarithmic depth. SIAM J. Discret. Math., 3(2):255–265, 1990. doi:10.1137/0403021.
\bibitem{r18} H. Kesten. An introduction to probability theory and its applications, volume i, (william feller). SIAM Review, 11(1):96–96, 1969. doi:10.1137/1011021.
\bibitem{r19} Erfan Khaniki. On proof complexity of resolution over polynomial calculus. Electronic Colloquium on Computational Complexity (ECCC), 27:34, 2020. URL: https://eccc.weizmann. ac.il/report/2020/034.
\bibitem{r20} Jan Krajíček. Interpolation theorems, lower bounds for proof systems, and independence results for bounded arithmetic. J. Symb. Log., 62(2):457–486, 1997. doi:10.2307/2275541.
\bibitem{r21} Jan Krajíček. An exponential lower bound for a constraint propagation proof system based on ordered binary decision diagrams. J. Symb. Log., 73(1):227–237, 2008. doi:10.2178/jsl/ 1208358751.
\bibitem{r22} Jan Krajíček. Randomized feasible interpolation and monotone circuits with a local oracle. J. Mathematical Logic, 18(2):1850012:1–1850012:27, 2018. doi:10.1142/S0219061318500125.
\bibitem{r23} Jan Krajíček. Proof complexity, volume 170. Cambridge University Press, 2019.
\bibitem{r24} Eyal Kushilevitz and Noam Nisan. Communication complexity. Cambridge University Press, 1997.
\bibitem{r25} Edward I Nechiporuk. A boolean function. Engl. transl. in Sov. Phys. Dokl., 10:591–593, 1966.
\bibitem{r26} Vsevolod Oparin. Tight upper bound on splitting by linear combinations for pigeonhole principle. In Nadia Creignou and Daniel Le Berre, editors, Theory and Applications of Satisfiability Testing - SAT 2016 - 19th International Conference, Bordeaux, France, July 5-8, 2016, Proceedings, volume 9710 of Lecture Notes in Computer Science, pages 77–84. Springer, 2016. doi:10.1007/978-3-319-40970-2\_6.
\bibitem{r27} Fedor Part and Iddo Tzameret. Resolution with counting: Dag-like lower bounds and different moduli. In Thomas Vidick, editor, 11th Innovations in Theoretical Computer Science Conference, ITCS 2020, January 12-14, 2020, Seattle, Washington, USA, volume 151 of LIPIcs, pages 19:1–19:37. Schloss Dagstuhl - Leibniz-Zentrum für Informatik, 2020. doi:10.4230/LIPIcs.ITCS.2020.19.
\bibitem{r28} Pavel Pudlák. Lower bounds for resolution and cutting plane proofs and monotone computations. J. Symb. Log., 62(3):981–998, 1997. doi:10.2307/2275583.
\bibitem{r29} Ran Raz and Pierre McKenzie. Separation of the monotone NC hierarchy. Combinatorica, 19(3):403–435, 1999. doi:10.1007/s004930050062.
\bibitem{r30} Ran Raz and Avi Wigderson. Monotone circuits for matching require linear depth. J. ACM, 39(3):736–744, 1992. doi:10.1145/146637.146684.
\bibitem{r31} Alexander A. Sherstov. The multiparty communication complexity of set disjointness. In Proceedings of the forty-fourth annual ACM symposium on Theory of computing, pages 525–548, 2012.
\bibitem{r32} Alexander A. Sherstov. Communication lower bounds using directional derivatives. J. ACM, 61(6), December 2014. doi:10.1145/2629334.
\bibitem{r33} Dmitry Sokolov. Dag-like communication and its applications. In Computer Science Theory and Applications - 12th International Computer Science Symposium in Russia, CSR 2017, Kazan, Russia, June 8-12, 2017, Proceedings, pages 294–307, 2017. doi:10.1007/ 978-3-319-58747-9\_26.
\end{thebibliography}
\end{document}
